\documentclass[11pt]{article}
\usepackage[a4paper,margin=26mm]{geometry}
\usepackage[T1]{fontenc}
\usepackage{lmodern,amsmath,amssymb,amsthm,mathtools,microtype}
\usepackage[colorlinks=true,linkcolor=blue,citecolor=blue,urlcolor=blue]{hyperref}
\newcommand{\D}{\mathbb D}
\newcommand{\E}{\mathbb E}
\newcommand{\Pp}{\mathbb P}
\newcommand{\dd}{\,d}
\DeclareMathOperator{\TV}{TV}
\DeclareMathOperator{\osc}{osc}
\newtheorem{theorem}{Theorem}[section]
\newtheorem{proposition}[theorem]{Proposition}
\newtheorem{lemma}[theorem]{Lemma}
\newtheorem{corollary}[theorem]{Corollary}
\theoremstyle{remark}\newtheorem{remark}[theorem]{Remark}
\title{Screening estimates for the actual unweighted disk gas\\[4pt]
\large Boundary attenuation and exactly centered reference gases}
\author{}
\date{8 September 2026}
\begin{document}
\maketitle
\vspace{-1.5em}
\begin{abstract}
We prove screening estimates beyond the determinantal case for the
original finite unweighted disk log gas. For every $0<\beta\le2$,
a negative-type distance argument gives a particle-number-uniform
susceptibility inequality under every common one-body weight. Its
consequences include generating-functional bounds and total-variation
stability of the complete particle law under perturbations with small
quadratic intensity cost. A positive radial-gap mixture supplies a
uniform hyperbolic intensity bound at every temperature. Combining
these facts, distant Green-potential tilts have vanishing effect on
the actual law, uniformly in $N$, and the full exterior Green force
has vanishing mean-square tail. The limiting Green force is therefore
recoverable from the limiting exterior configuration. Independently,
the nondegenerate outermost-radius boundary statistic decouples from
the entire local law at rate $O(N^{-1})$ for every $\beta>0$.
At $\beta=2$ we prove full selected-exterior attenuation at rate
$O(1-R^2)$, uniformly in $N$, by reversing the conditioning and
computing the distant response to an interior hole.
We also construct an exact short-range reference representation,
including its count weights, with a uniformly small comparison norm
for $0<\beta\le1/4$. For $0<\beta\le1/20$, a convergent complex
cluster expansion gives a periodic restoring potential in the exact
imaginary Gaussian-field representation. A self-consistent reference
then cancels the linear field term exactly while retaining these
uniform small-temperature bounds. A separate strongly convex
construction gives exact canonical centering at every $N$ and
$0<\beta<2$. The remaining large-field and normalization problems
are not discarded.
Appendices give a non-determinantal two-to-one deletion coupling and
whole-local-law linear response at $\beta=2$. Independence from the
entire distant exterior away from $2$, and the conjectured intensity
coefficient, remain unproved.
\end{abstract}
\clearpage
\tableofcontents
\clearpage

\section{The actual model and the scope of the results}
The finite gas throughout the note is
\[
 P_{N,\beta}(d\mathbf z)
 =\frac1{Z_N(\beta)}\prod_{i<j}|z_i-z_j|^\beta
                   \prod_i\mu(dz_i),\qquad
 \mu=\frac{dA}{\pi},\quad z_i\in\D,\quad\beta>0.
\]
Write $\Xi_N=\sum_i\delta_{z_i}$ and $B_r=\{|z|<r\}$.
Intensities $\rho_{N,\beta}$ are relative to $\mu$. Use curvature
$-1$ hyperbolic distance $d$ and area
\[
 \nu(dz)=\frac{\mu(dz)}{(1-|z|^2)^2},\qquad
 d(0,z)=2\operatorname{artanh}|z|.
\]
The hyperbolic ball of radius $L$ has $\nu$-volume $\sinh^2(L/2)$.
For probabilities, $\TV(P,Q)=\sup_E|P(E)-Q(E)|$ takes values in
$[0,1]$. Thus a TV estimate for a point-process restriction controls
all its positions and count probabilities at once.

The full boundary target is, for fixed $r<1$,
\begin{equation}\label{eq:full-exterior-target}
 \sup_N\mathbb E\TV\left(
  \mathcal L(\Xi_N|_{B_r}\mid\Xi_N|_{\D\setminus B_R}),
  \mathcal L(\Xi_N|_{B_r})\right)\longrightarrow0
       \quad(R\uparrow1).
\end{equation}
The results below establish specific parts of this program at fixed
temperatures. Theorem~\ref{thm:full-exterior-beta-two} proves
\eqref{eq:full-exterior-target} at $\beta=2$, including all local
counts, for boundary data selected by the actual finite gases.
It remains unproved away from the determinantal point.
The following distinctions are maintained:
\begin{enumerate}
\item The susceptibility and Green-tail estimates use the original
canonical gas and hold for the whole interval $0<\beta\le2$.
\item The radial-scale theorem conditions on one nondegenerate
statistic measurable from a distant exterior; it does not condition
on every detail of that exterior.
\item The small reference norm belongs to a specified short-range
reference law. The exact original law retains an additional tilt.
\item Exact self-consistent centering removes the reference's linear
field term. The canonical construction does not automatically
inherit the grandcanonical influence bound or its phase coercivity.
\item The temperature-response theorem concerns derivatives taken
at $2$ before $N\to\infty$; it does not control a fixed nonlinear
temperature interval.
\end{enumerate}

The central new input is
\begin{equation}\label{eq:intro-susceptibility}
 \operatorname{Var}\Xi_N(f)
 \le\mathbb E\Xi_N(f^2)-\frac{(\mathbb E\Xi_N(f))^2}{N},
 \qquad 0<\beta\le2.
\end{equation}
It remains valid under arbitrary common finite one-body weights,
including exact finite exterior weights. The proof below is a
self-contained Hilbert-space and spectral induction. Its connection
with Lorentzian polynomials uses the original theorem of
Br\"and\'en--Huh~\cite{BrandenHuh}; that theorem is not needed for
the direct continuum proof. This supplies a new input to the
screening program studied in the preceding note, rather than simply
repeating a comparison criterion with unverified hypotheses.

All asserted estimates are proved in the following sections. Standard
spectral, measure-theoretic, elliptic, and determinantal identities
used in the arguments are specified there. The longest independent
mathematical checks concerned the spectral induction, the radial
transport, the Green-field limit, the exact reference normalization,
the all-order signed Janossy inversion, the full-exterior comparison,
and the new self-consistent centering constructions.


\section{A count-sensitive susceptibility bound for the actual gas below two}
\label{sec:lorentzian}

This section gives an estimate for the original interaction, uniformly in
particle number and in every finite exterior weight.  It uses the fact
that $|z-w|^\beta$ is a squared Hilbert distance for $0<\beta\le2$.
The corresponding finite generating polynomials are Lorentzian.  We
also give a direct spectral induction proof of the needed consequence,
so the continuum estimate does not depend on a discretization argument.

\begin{lemma}[A Hilbert representation]
\label{lem:distance-negative-type}
For $0<\beta\le2$ there are a real Hilbert space $H$ and a continuous
map $\Phi:\mathbb R^2\to H$ such that
\[
 \|\Phi(x)-\Phi(y)\|_H^2=|x-y|^\beta.
\]
Consequently, for any finite points $x_1,\ldots,x_q$ and real numbers
$c_i$ with $\sum_i c_i=0$,
\[
 \sum_{i,j}c_ic_j|x_i-x_j|^\beta\le0.
\]
\end{lemma}

\begin{proof}
For $\beta=2$, take $H=\mathbb R^2$ and $\Phi(x)=x$.
For $0<\beta<2$, put
\[
 C_\beta=\int_{\mathbb R^2}
       |e^{ik\cdot e_1}-1|^2|k|^{-2-\beta}\,dk.
\]
This number is finite and strictly positive: near zero the numerator
is at most $|k|^2$, and outside the unit ball it is at most four.
Use the underlying real Hilbert space of
$L^2(\mathbb R^2,|k|^{-2-\beta}dk;\mathbb C)$, and set
$\Phi(x)(k)=C_\beta^{-1/2}(e^{ik\cdot x}-1)$.
Rotation invariance and the substitution $u=|x-y|k$ give the displayed
distance identity.  That identity also proves continuity.  Finally,
expanding squared Hilbert distances, with $\sum_i c_i=0$, gives
\[
 \sum_{i,j}c_ic_j\|\Phi(x_i)-\Phi(x_j)\|^2
 =-2\Big\|\sum_i c_i\Phi(x_i)\Big\|^2\le0.
\]
\end{proof}

\begin{proposition}[Lorentzian finite approximations]
\label{prop:finite-lorentzian}
Let $x_1,\ldots,x_q$ be distinct planar points, let $w_i>0$, and let
$1\le n\le q$.  If $0<\beta\le2$, then the polynomial
\[
 F_n(t_1,\ldots,t_q)=
 \sum_{\substack{S\subset\{1,\ldots,q\}\\|S|=n}}
  \left(\prod_{\{i,j\}\subset S}|x_i-x_j|^\beta\right)
  \prod_{i\in S}w_it_i
\]
is Lorentzian.  In particular, $\log F_n$ and $F_n^{1/n}$ are concave
on the positive orthant.
\end{proposition}

\begin{proof}
We use the characterization and concavity theorem of
Br\"and\'en--Huh \cite[Theorems 2.25 and 2.30, Proposition 2.33]{BrandenHuh}:
a homogeneous polynomial with nonnegative coefficients and
$M$-convex support is Lorentzian if every derivative of degree two
has a Hessian with at most one positive eigenvalue; such a polynomial
has the two stated concavity properties.

Here the support consists of all zero-one vectors with coordinate sum
$n$.  It satisfies the required exchange property: if $i\in S\setminus
T$, choose any $j\in T\setminus S$ and exchange $i,j$.
For $n\ge2$, a nonzero derivative indexed by an $(n-2)$-element set
$A$ has quadratic coefficients
\[
 C_A\,b_i b_j |x_i-x_j|^\beta,
 \qquad
 b_i=w_i\prod_{a\in A}|x_i-x_a|^\beta,
 \qquad i,j\notin A,
\]
with $C_A>0$.  Its Hessian is thus a positive diagonal congruence of
the matrix $K_{ij}=|x_i-x_j|^\beta$, with zero diagonal.
Lemma~\ref{lem:distance-negative-type} makes the quadratic form of $K$
nonpositive on the codimension-one subspace $\sum_i c_i=0$.
The min--max principle implies that $K$ has at most one positive
eigenvalue, and diagonal congruence preserves its inertia.  Derivatives
with repeated indices vanish.  This checks every hypothesis.  Degree
one is immediate.
\end{proof}

\begin{theorem}[Uniform susceptibility for every finite exterior weight]
\label{thm:subpoisson-susceptibility}
Let $\sigma$ be a nonzero finite nonatomic Borel measure with compact
support in $\mathbb R^2$, let $n\ge1$, and let $0<\beta\le2$.
Normalize the symmetric probability measure
\[
 P_{n,\beta}^{\sigma}(d\mathbf z)
 =\frac{1}{Z_{n,\beta}^{\sigma}}
   \prod_{i<j}|z_i-z_j|^\beta\prod_{i=1}^n\sigma(dz_i).
\]
Write $\Xi=\sum_i\delta_{z_i}$ and let $\rho$ be its intensity measure,
so $\rho(\mathbb R^2)=n$.  Every bounded real Borel function $f$ satisfies
\begin{equation}
 \operatorname{Var}\Xi(f)
 \le \int f^2\,d\rho-\frac1n\left(\int f\,d\rho\right)^2
 \le \int f^2\,d\rho.
 \label{eq:subpoisson-variance}
\end{equation}
For every bounded nonnegative Borel function $v$,
\begin{equation}
 \mathbb E\prod_{z\in\Xi}v(z)
 \le\left(\frac1n\int v\,d\rho\right)^n
 \le\exp\!\left\{\int(v-1)\,d\rho\right\}.
 \label{eq:subpoisson-generating}
\end{equation}
Consequently, for every real $t$,
\begin{equation}
 \log\mathbb E e^{t(\Xi(f)-\mathbb E\Xi(f))}
 \le\int(e^{tf}-1-tf)\,d\rho.
 \label{eq:subpoisson-laplace}
\end{equation}
These inequalities hold with the same constants after replacing
$\sigma$ by any nonzero finite nonatomic one-body weight.
\end{theorem}

\begin{proof}[Direct proof of the variance estimate]
Compact support makes the interaction bounded and the partition
function finite.  Nonatomicity makes the partition function positive,
since the particles are pairwise distinct almost everywhere.
Normalize $\sigma$ to have mass one, which does not change the law.
We prove the first inequality by induction on $n$.  For $n=1$ it is
an equality.

For $n=2$, use Lemma~\ref{lem:distance-negative-type} and write
$m=\int\Phi\,d\sigma$, $H(z)=\Phi(z)-m$, and
$V=\int\|H\|^2d\sigma>0$.  The pair partition function is $2V$.
Let $f$ have mean zero under the tagged-particle marginal, and put
\[
 a=\int f\,d\sigma,\qquad
 b=\int fH\,d\sigma,\qquad
 c=\int f\|H\|^2d\sigma.
\]
The tagged marginal has density $(\|H(z)\|^2+V)/(2V)$ relative to
$\sigma$, so the centering condition is $c+Va=0$.  Expanding the
Hilbert distance gives
\[
 \mathbb E[f(z_1)f(z_2)]
 =\frac{2ac-2\|b\|^2}{2V}
 =-a^2-\frac{\|b\|^2}{V}\le0.
\]
Thus $\operatorname{Var}(f(z_1)+f(z_2))\le2\mathbb E f(z_1)^2$,
as required for centered $f$.

Suppose now that the assertion holds for $n-1\ge2$, for every such
one-body measure.  Let $\pi$ be the tagged-particle marginal of the
$n$-particle law, and define the self-adjoint Markov operator on
$L^2(\pi)$ by
\[
 (Tf)(x)=\mathbb E[f(z_2)\mid z_1=x].
\]
Exchangeability gives self-adjointness, and conditional Jensen gives
$\|T\|\le1$.  Take bounded $f$ with $\pi(f)=0$, and write
\[
 M=\pi(f^2),\qquad A=\langle f,Tf\rangle_\pi,
 \qquad B=\|Tf\|_{L^2(\pi)}^2.
\]
Conditioned on $z_1=x$, the other $n-1$ particles have precisely the
same interaction and one-body measure $|x-z|^\beta\sigma(dz)$.
This measure remains finite and nonatomic.  The induction hypothesis
and total variance therefore give
\begin{align*}
 \operatorname{Var}\sum_{i=1}^nf(z_i)
 &\le (n-1)(M-B)
       +\operatorname{Var}\bigl(f(z_1)+(n-1)Tf(z_1)\bigr)\\
 &=nM+2(n-1)A+(n-1)(n-2)B.
\end{align*}
On the other hand, exchangeability gives the exact expression
$nM+n(n-1)A$ for the left side.  Hence $A\le B$, since $n>2$.
Bounded functions are dense in $L^2(\pi)$, so on the mean-zero
subspace this is the operator inequality $T\le T^2$.
The spectral theorem and $\|T\|\le1$ place its spectrum in
$[-1,0]\cup\{1\}$.

The eigenvalue one has only constant eigenfunctions.  Indeed,
\[
 \langle f,(I-T)f\rangle_\pi
 =\tfrac12\mathbb E\bigl(f(z_1)-f(z_2)\bigr)^2,
\]
and the pair marginal has a strictly positive density relative to
$\sigma\otimes\sigma$ almost everywhere.  To see positivity, fix two
distinct points; the integral over the other $n-2$ coordinates is
strictly positive because nonatomicity excludes all collisions almost
everywhere.  Thus equality in the displayed formula forces $f$ to
be constant $\sigma$-almost everywhere, hence also $\pi$-almost
everywhere.  Since one is isolated from the remaining possible
spectrum, its spectral subspace is exactly this eigenspace.
There is therefore no positive spectrum on the mean-zero subspace,
and $A\le0$.  It follows that
\[
 \operatorname{Var}\sum_i f(z_i)
 =nM+n(n-1)A\le nM.
\]
Subtracting the tagged mean from an arbitrary bounded $f$ proves
\eqref{eq:subpoisson-variance} and completes the induction.
\end{proof}

\begin{proof}[Generating-functional and Laplace bounds]
First suppose that $v$ is bounded and bounded away from zero.  Set
$v_s=1+s(v-1)$, $0\le s\le1$, and
\[
 G(s)=\mathbb E\prod_{z\in\Xi}v_s(z),\qquad L(s)=\log G(s).
\]
All derivatives can be taken under the finite integral.  Under the
tilted gas with one-body measure $v_s\sigma$, put
$h_s=(v-1)/v_s$.  Exact differentiation gives
\[
 L'(s)=\mathbb E_s\Xi(h_s),\qquad
 L''(s)=\operatorname{Var}_s\Xi(h_s)
                -\mathbb E_s\Xi(h_s^2)
 \le-\frac{L'(s)^2}{n},
\]
where the last step is \eqref{eq:subpoisson-variance} for that tilted
one-body measure.  Consequently $(G^{1/n})''\le0$.
Its tangent line at zero gives
\[
 G(1)^{1/n}\le1+\frac1n\int(v-1)\,d\rho
 =\frac1n\int v\,d\rho.
\]
Apply this first to $v+\varepsilon$ and let $\varepsilon\downarrow0$
to include bounded nonnegative $v$.  The exponential inequality
follows from $(1+u/n)^n\le e^u$ for $u\ge-n$, with the zero case
interpreted directly.  Finally set $v=e^{tf}$ and subtract
$t\int f\,d\rho$ from the logarithm to obtain
\eqref{eq:subpoisson-laplace}.
\end{proof}

\begin{corollary}[Count sectors of the actual conditional gas]
\label{cor:conditional-count-susceptibility}
Let $U$ be any bounded observation region, let $Y$ be a finite exterior
configuration, and condition the disk gas on having exactly $m$ points
in $U$, with exterior $Y$.  For $0<\beta\le2$, all conclusions of
Theorem~\ref{thm:subpoisson-susceptibility} hold with $n=m$ and
\[
 \sigma(dz)=\mathbf1_U(z)\prod_{y\in Y}|z-y|^\beta\,dA(z).
\]
The empty sector is deterministic.  In any nonempty sector, for a
Borel set $A\subset U$ and $a=\mathbb E\Xi(A)$,
\[
 \operatorname{Var}\Xi(A)\le a(1-a/m),\qquad
 \mathbb E t^{\Xi(A)}
 \le\left(1+\frac am(t-1)\right)^m\le e^{a(t-1)},\quad t>0.
\]
\end{corollary}

\begin{proof}
The displayed measure is finite and nonatomic, and nonzero whenever
the sector is admissible.  Apply the theorem with $f=\mathbf1_A$ and
$v=1+(t-1)\mathbf1_A$.  This proof requires no separation of $Y$ into
annuli and no approximation of its entire interaction.
\end{proof}

\begin{remark}[Scope of the negative dependence]
The theorem controls the variance of every real linear statistic and
the full generating functional, including count fluctuations.  It
does not assert that the covariance of two different nonnegative
test functions is always nonpositive, nor that arbitrary local events
are negatively associated.  It does not by itself provide stochastic
domination between successive particle numbers.  A spatial attenuation
estimate still needs a perturbation whose quadratic intensity cost
vanishes with distance.
\end{remark}


\section{Radial comparison and an actual screened boundary statistic}
\label{sec:radial-comparison}

All statements in this section concern the unweighted disk gas at an
arbitrary fixed $\beta>0$. Intensities are relative to
$\mu(\dd z)=\dd A(z)/\pi$. Set
\[
 d_j=2j+\frac\beta2j(j-1),\qquad
 c_\beta=\min\{2,\beta/2\},\qquad
 C_\beta=2+\frac{16\sqrt6}{c_\beta}.
\]
The notation $c_\beta$ here is merely a comparison constant, not the
conjectured hyperbolic intensity coefficient.

\begin{proposition}[A positive mixture for the ordered radial gaps]
\label{prop:radial-positive-mixture}
Order the radii increasingly as $r_1<\cdots<r_N$, put
$t_j=-\log r_j$, $t_{N+1}=0$, and $g_j=t_j-t_{j+1}$.
The joint distribution of $(g_1,\ldots,g_N)$ is a countable mixture of
product distributions
\[
 \bigotimes_{j=1}^N\operatorname{Exp}(d_j+2A_j),
 \qquad A_j\in\mathbb Z_{\ge0},\quad A_N=0.
\]
In particular, on a coupling with independent rate-one exponentials
$E_j$, one has simultaneously
\begin{equation}\label{eq:radial-gap-domination}
 g_j\le E_j/d_j,\qquad
 t_j\le\sum_{k=j}^N E_k/d_k.
\end{equation}
\end{proposition}

\begin{proof}
Ties and zero radii have probability zero. On the ordered region factor
\[
 \prod_{i<j}|z_j-z_i|^\beta
 =\prod_{j=1}^N r_j^{\beta(j-1)}
   \left|\prod_{i<j}(1-z_i/z_j)^{\beta/2}\right|^2.
\]
Every power is the analytic branch defined by its binomial expansion
on $|z_i/z_j|<1$. The finite product has an absolutely convergent
Laurent expansion on every fixed ordered radial torus,
\[
 \prod_{i<j}(1-z_i/z_j)^{\beta/2}
   =\sum_\alpha b_\alpha\prod_j z_j^{\alpha_j}.
\]
An expansion term is specified by nonnegative integers $n_{ij}$,
$i<j$, and has
\[
 \alpha_j=\sum_{\ell>j}n_{j\ell}-\sum_{i<j}n_{ij},
 \qquad
 A_k:=\sum_{j\le k}\alpha_j
       =\sum_{i\le k<j}n_{ij}\ge0,
 \qquad A_N=0.
\]
For a fixed exponent vector, each $n_{ij}$ is bounded by any $A_k$
with $i\le k<j$, so only finitely many terms contribute to its
coefficient. Absolute convergence also justifies regrouping terms.
Angular Parseval therefore gives
\[
 \int\left|\prod_{i<j}(1-z_i/z_j)^{\beta/2}\right|^2
       \prod_j\frac{\dd\theta_j}{2\pi}
 =\sum_\alpha |b_\alpha|^2\prod_j r_j^{2\alpha_j}.
\]
The radial area elements are $2r_j\dd r_j$. After the changes to
$t_j$ and then $g_j$, the contribution of a term is, up to a common
positive constant,
\[
 |b_\alpha|^2
 \exp\left\{-\sum_{k=1}^N(d_k+2A_k)g_k\right\}
 \prod_k\dd g_k.
\]
Indeed, the coefficient of $g_k$ is
$\sum_{j\le k}[2+\beta(j-1)+2\alpha_j]=d_k+2A_k$.
All rates are positive. Tonelli's theorem applies to this nonnegative
series. Its integral is finite and nonzero because it is the original
finite partition integral. Normalizing its terms gives mixture weights
proportional to
$|b_\alpha|^2\prod_k(d_k+2A_k)^{-1}$ and proves the claim.
Finally, sample the mixture index and independent $E_j$, and set
$g_j=E_j/(d_j+2A_j)$ to obtain \eqref{eq:radial-gap-domination}.
\end{proof}

\begin{theorem}[Uniform compact intensity bound]
\label{thm:uniform-actual-intensity}
For all $N\ge1$, $\beta>0$, and $z\in\D$,
\begin{equation}\label{eq:actual-uniform-bergman-upper}
 \rho_{N,\beta}(z)\le\frac{C_\beta}{(1-|z|^2)^2}.
\end{equation}
No determinantal representation is used.
\end{theorem}

\begin{proof}
The inequality $d_k\ge c_\beta k^2$ follows directly from its
definition. Independence of the exponentials in
\eqref{eq:radial-gap-domination} gives
\[
 \E t_j^2
 \le\left(\sum_{k=j}^N d_k^{-1}\right)^2
       +\sum_{k=j}^N d_k^{-2}
 \le\frac{6}{c_\beta^2j^2}.
\]
Here $\sum_{k\ge j}k^{-2}\le2/j$ and
$\sum_{k\ge j}k^{-4}\le2/j^2$ suffice. Since
$\{\Xi_N(B_r)\ge j\}=\{t_j> -\log r\}$ up to null events,
Markov's inequality yields, for $0<r<1$,
\begin{equation}\label{eq:actual-radial-count-upper}
 \E\Xi_N(B_r)
 \le\sum_{j\ge1}\min\left\{1,
       \frac{6}{c_\beta^2j^2(-\log r)^2}\right\}
 \le1+\frac{2\sqrt6}{c_\beta(-\log r)}.
\end{equation}
For the last inequality, put $a=\sqrt6/(c_\beta(-\log r))$,
split the sum at $\lceil a\rceil$, and use
$\sum_{j>m}j^{-2}\le1/m$.

The intensity is rotation invariant and subharmonic in its argument:
its marginal formula is a positive integral of
$|\prod_{i=2}^N(z-z_i)|^\beta$ against a measure independent of $z$.
These integrands are subharmonic, and are locally uniformly bounded
for fixed $N$. A radial subharmonic function finite at zero is
nondecreasing with radius, as follows from the monotonicity of its
circular means. Consequently, whenever $0\le r<s<1$,
\[
 \rho_{N,\beta}(r)(s^2-r^2)
 \le\int_{B_s\setminus B_r}\rho_{N,\beta}\,\dd\mu
 \le\E\Xi_N(B_s).
\]
Choose $s^2=(1+r^2)/2$. Then
$-\log s\ge(1-s^2)/2=(1-r^2)/4$. Substitution of
\eqref{eq:actual-radial-count-upper} gives
\[
 \rho_{N,\beta}(r)
 \le\frac2{1-r^2}
       +\frac{16\sqrt6}{c_\beta(1-r^2)^2}
 \le\frac{C_\beta}{(1-r^2)^2}.
\]
\end{proof}

\begin{proposition}[Exact scale separation]
\label{prop:exact-scale-separation}
Let $R_N=\max_i|z_i|$, $W_N=\{z_i/R_N:1\le i\le N\}$, and
$T_N=d_N(-\log R_N)$. Then $T_N$ has the rate-one exponential
law and is independent of the whole normalized configuration $W_N$.
The intensity of $W_N$ in the open disk satisfies the same bound
\eqref{eq:actual-uniform-bergman-upper}.
\end{proposition}

\begin{proof}
For completeness, independence of the whole shape, rather than only
its radii, follows from a direct change of variables. Choose the
almost surely unique outermost label, write its coordinate as
$R e^{i\theta}$, and write all other coordinates as $R w_i$,
$|w_i|<1$. The Jacobian and the homogeneous Vandermonde yield a
factor $R^{d_N-1}\dd R$. All remaining factors depend only on
$(\theta,w_1,\ldots,w_{N-1})$. Thus $R$ has density
$d_N R^{d_N-1}$ on $(0,1)$ independently of the shape, proving
the first assertion. Equivalently, the gap mixture has $A_N=0$.

There is one shape particle on the unit circle. The other particles,
conditional on its angle, have density proportional to
\[
 |\Delta_{N-1}(\mathbf w)|^\beta
 \prod_{i=1}^{N-1}|w_i-e^{i\theta}|^\beta
 \prod_i\mu(\dd w_i).
\]
The interior shape intensity, after averaging the angle, is therefore
radial and subharmonic by exactly the marginal argument above.
Moreover $R_N\le1$ gives, on the defining coupling,
$W_N(B_s)\le\Xi_N(B_s)$ for every $s<1$. Hence
\eqref{eq:actual-radial-count-upper} applies to its count expectation.
The same monotonicity argument proves its stated intensity bound.
For $N=1$ its interior intensity is zero, so the conclusion still holds.
\end{proof}

\begin{theorem}[Count-law attenuation of the outermost radial mode]
\label{thm:radial-count-screening}
Fix $0<r<1$, let $s^2=(1+r^2)/2$, and put
$K_{\beta,r}=C_\beta/(1-s^2)^2$.
For $N\ge2$, with $M_N=\Xi_N(B_r)$, one has
\begin{equation}\label{eq:radial-count-TV}
 \E\TV\bigl(\mathcal L(M_N\mid T_N),\mathcal L(M_N)\bigr)
 \le\frac{4K_{\beta,r}r^2}{d_N-2}
           +2(r/s)^{d_N}
 =O_{\beta,r}(N^{-2}).
\end{equation}
This is an actual selected-boundary estimate for every fixed
$\beta>0$ and includes the complete probability distribution of
the local count. The boundary statistic $T_N$ remains nondegenerate.
\end{theorem}

\begin{proof}
Let $Q_t$ be the law of $W_N(B_{r e^{t/d_N}})$, with the evident
interpretation if the radius exceeds one, and let $Q_0$ be the law
at radius $r$. Exact scale separation says that $Q_t$ is a version
of the conditional law of $M_N$ given $T_N=t$. Couple these count
variables using the same $W_N$. If
$0\le t\le d_N\log(s/r)$, the coupling inequality and the shape
intensity bound give
\[
 \TV(Q_t,Q_0)
 \le \Pp\{W_N(B_{r e^{t/d_N}}\setminus B_r)>0\}
 \le K_{\beta,r}r^2(e^{2t/d_N}-1).
\]
For larger $t$, use the bound one. Since $d_N>2$ for $N\ge2$,
integration against the exponential density yields
\[
 \int_0^\infty\TV(Q_t,Q_0)e^{-t}\dd t
 \le\frac{2K_{\beta,r}r^2}{d_N-2}+(r/s)^{d_N}.
\]
The unconditional law is $\overline Q=\int Q_t e^{-t}\dd t$.
Convexity and the triangle inequality imply
$\int\TV(Q_t,\overline Q)e^{-t}\dd t
 \le2\int\TV(Q_t,Q_0)e^{-t}\dd t$,
proving the claim.
\end{proof}

The same proof controls the joint vector of counts in any fixed finite
collection of concentric disks: replace $r^2$ in the first term by
the sum of their squared radii, use their largest radius in the tail
term, and choose $s$ above that largest radius. Counts in finitely many
concentric annuli follow by taking differences. This does not show
independence from the \emph{entire} distant exterior, whose angular
information and remaining count selection are absent from $T_N$.
The result is useful because it separates a rigorously attenuated,
nontrivial boundary mode from those still requiring collective screening.

\begin{lemma}[Uniform exponential compact-count moments]
\label{lem:radial-exponential-count}
For $0<b<1$ and every integer $j\ge1$,
\begin{equation}\label{eq:radial-Gaussian-count-tail}
 \Pp\{\Xi_N(B_b)\ge j\},\quad
 \Pp\{W_N(B_b)\ge j\}
 \le\exp\left\{2j-
           \frac{c_\beta(-\log b)}2j^2\right\}.
\end{equation}
In particular, for every fixed $v>0$,
\[
 \sup_N\E\bigl[\Xi_N(B_b)e^{v\Xi_N(B_b)}\bigr]<\infty,
 \qquad
 \sup_N\E\bigl[W_N(B_b)e^{vW_N(B_b)}\bigr]<\infty.
\]
\end{lemma}

\begin{proof}
For $j\le N$ put $q=c_\beta j^2/2$. Then
$q/d_k\le j^2/(2k^2)\le1/2$ for $k\ge j$. The gap comparison and
$-\log(1-x)\le2x$ on $[0,1/2]$ yield
\[
 \log\E\exp\left(q\sum_{k=j}^N E_k/d_k\right)
 \le2q\sum_{k=j}^N d_k^{-1}\le2j.
\]
Exponential Markov proves the first tail bound; it is automatic for
$j>N$. The shape bound follows from its coupled count domination.
For either count $K$, monotone summation gives
\[
 \E[K e^{vK}]
 =\sum_{j\ge1}\bigl(je^{vj}-(j-1)e^{v(j-1)}\bigr)\Pp(K\ge j)
 \le\sum_{j\ge1}j\exp\left\{(v+2)j-
               \frac{c_\beta(-\log b)}2j^2\right\}<\infty.
\]
\end{proof}

\begin{theorem}[Full local-law attenuation of the radial scale mode]
\label{thm:radial-full-law-screening}
For every fixed $\beta>0$ and $0<r<1$, there are constants
$A_{\beta,r}<\infty$ and $a_{\beta,r}>0$, independent of $N$, such that
\begin{equation}\label{eq:radial-full-law-TV}
 \E\TV\left(
    \mathcal L(\Xi_N|_{B_r}\mid T_N),
    \mathcal L(\Xi_N|_{B_r})\right)
 \le A_{\beta,r}\frac{N}{d_N}
             +2\exp\{-a_{\beta,r}d_N/N\}
 =O_{\beta,r}(N^{-1}).
\end{equation}
This concerns the complete local law, including all particle positions
and particle-number probabilities.
\end{theorem}

\begin{proof}
Choose $r<u<b<1$ and a smooth radial function $\chi$ equal to one
on $B_u$ and zero outside $B_b$. Put $V(z)=\chi(z)z$ and
$L=\max\{1,\operatorname{Lip}(V)\}$. For
\[
 0\le\eta\le\eta_0
 :=\min\left\{\frac1{2L},\frac{u-r}{2u},\frac12\right\},
 \qquad F_\eta(z)=z-\eta V(z),
\]
the map $F_\eta$ is a smooth diffeomorphism of the plane, equal to
the identity outside $B_b$. Indeed, its derivative has singular values
between $1-\eta L$ and $1+\eta L$, the map is globally injective by
the same Lipschitz bounds, and it is proper; local invertibility and
properness give surjectivity. It consequently preserves $B_b$ and
$\D$. Its pair distortions and Jacobian satisfy
\begin{align*}
 \left|\log\frac{|F_\eta(x)-F_\eta(y)|}{|x-y|}\right|
    &\le2L\eta \quad (x\ne y),\\
 |\log\det DF_\eta(x)|&\le4L\eta.
\end{align*}
The determinant is positive, by continuity from $\eta=0$. Pair
distortions vanish when both points lie outside $B_b$, and the
Jacobian is one there.

Use the shape coordinates from Proposition~\ref{prop:exact-scale-separation}:
one point is $e^{i\theta}$ on the unit circle and the remaining
$N-1$ points have a density $p_\theta(\mathbf w)$. The transport
fixes the boundary point. Write $K=W_N(B_b)$ and
$B=L(2\beta+4)$. There are at most $NK$ affected pairs and at most
$K$ nontrivial one-particle Jacobians. Therefore, almost surely,
\[
 H_\eta:=\log\frac{p_\theta(F_\eta\mathbf w)
                         J_{F_\eta}(\mathbf w)}
                        {p_\theta(\mathbf w)}
 \quad\hbox{satisfies}\quad
 |H_\eta|\le B\eta NK.
\]
The densities vanish only on collision sets of measure zero, which do
not affect this formula. Averaging the usual change-of-variables
identity for total variation over $\theta$ gives
\[
 \TV\bigl(\mathcal L(F_\eta W_N),\mathcal L(W_N)\bigr)
 \le\tfrac12\E|e^{H_\eta}-1|
 \le\tfrac12 B\eta N\E[K e^{B\eta NK}].
\]
This inequality first holds for the labelled shape coordinates and
then for the unordered process by data processing. By
Lemma~\ref{lem:radial-exponential-count}, the constant
$D=\sup_N\E[K e^K]$ is finite. Whenever $B\eta N\le1$, the last
display is at most $BD\eta N/2$.

Because $F_\eta$ is radial, injective, and equals multiplication by
$1-\eta$ on $B_u$, while $(1-\eta)u>r$, its restriction satisfies
the exact configuration identity
\[
 (F_\eta W_N)|_{B_r}=((1-\eta)W_N)|_{B_r}.
\]
Let $Q_t$ now denote the law of
$(e^{-t/d_N}W_N)|_{B_r}$ and let $Q_0$ be the local shape law.
Set $a=\min\{\eta_0,1/B\}$. For $0\le t\le a d_N/N$,
$\eta=1-e^{-t/d_N}\le t/d_N$ satisfies both required smallness
conditions, and hence
\[
 \TV(Q_t,Q_0)\le\frac{BDN}{2d_N}t.
\]
For larger $t$ use one. Thus
\[
 \int_0^\infty\TV(Q_t,Q_0)e^{-t}\dd t
 \le\frac{BDN}{2d_N}+e^{-a d_N/N}.
\]
Exact scale separation identifies $Q_t$ as the conditional local law
given $T_N=t$. Comparing with its exponential mixture and using the
same triangle/convexity argument as in
Theorem~\ref{thm:radial-count-screening} proves
\eqref{eq:radial-full-law-TV}, with $A_{\beta,r}=BD$ and
$a_{\beta,r}=a$.
\end{proof}

This really is a boundary statistic. For any fixed $r<R<1$, define
$T_{N,R}=T_N$ if $R_N>R$ and set $T_{N,R}=+\infty$ otherwise.
It is measurable from the exterior configuration in $\D\setminus B_R$
and is a deterministic function of $T_N$. Data processing therefore
preserves both attenuation bounds when $T_N$ is replaced by $T_{N,R}$.
The censoring probability is exactly $R^{d_N}$, so this selected
exterior statistic has an asymptotically nondegenerate exponential
distribution. The theorem does not assert that it is a sufficient
statistic for the entire exterior.


\section{Fixed-temperature screening estimates for the actual gas}
\label{sec:actual-screening}

We now combine the susceptibility estimate with the radial bound.
Every statement in this section uses the original canonical disk law;
no effective pair potential replaces its interaction. Write
$\nu(dz)=dA(z)/[\pi(1-|z|^2)^2]$ and retain the intensity constant
$C_\beta$ from Theorem~\ref{thm:uniform-actual-intensity}.

\begin{theorem}[Full-law stability under a one-body perturbation]
\label{thm:actual-tilt-TV}
Let $0<\beta\le2$, let $f$ be bounded and real, and put
\[
 \frac{dP_{N,\beta}^{t,f}}{dP_{N,\beta}}
 =\frac{e^{t\Xi_N(f)}}{\mathbb E_{N,\beta}e^{t\Xi_N(f)}},
 \qquad M=\|f\|_\infty,\quad
 Q_0=\mathbb E_{N,\beta}\Xi_N(f^2).
\]
For every real $t$,
\begin{align}
 D_{\rm KL}(P_{N,\beta}^{t,f}\Vert P_{N,\beta})
 &\le\frac{t^2}{2}e^{|t|M}Q_0,
 \label{eq:actual-tilt-KL}\\
 \TV(P_{N,\beta}^{t,f},P_{N,\beta})
 &\le\frac{|t|}{2}e^{|t|M/2}\sqrt{Q_0}
 \le\frac{|t|\sqrt{C_\beta}}2 e^{|t|M/2}\|f\|_{L^2(\nu)}.
 \label{eq:actual-tilt-TV}
\end{align}
The last inequality is useful when $f\in L^2(\nu)$. The first two
apply more generally to every nonzero finite nonatomic one-body measure
supported in a bounded region, with
its own value of $Q_0$. The TV estimate concerns the entire
$N$-particle configuration, and therefore every local count and
position event simultaneously.
\end{theorem}

\begin{proof}
Let $P_s=P_{N,\beta}^{s,f}$ and
$Q(s)=\mathbb E_s\Xi_N(f^2)$. Differentiation is justified because
$\Xi_N(f)$ is bounded for fixed $N$. Theorem~
\ref{thm:subpoisson-susceptibility} applies to every $P_s$. Thus
\[
 |Q'(s)|
 =|\operatorname{Cov}_s(\Xi_N(f^2),\Xi_N(f))|
 \le\sqrt{\mathbb E_s\Xi_N(f^4)\,
                     \mathbb E_s\Xi_N(f^2)}
 \le M Q(s).
\]
Gronwall's inequality gives $Q(s)\le e^{|s|M}Q(0)$.
For $\psi(s)=\log\mathbb E_0e^{s\Xi_N(f)}$, this implies
$\psi''(s)=\operatorname{Var}_s\Xi_N(f)\le e^{|s|M}Q_0$.
The relative entropy is $t\psi'(t)-\psi(t)$, since $\psi(0)=0$.
Its derivative is $t\psi''(t)$; integration from zero to $t$,
also when $t<0$, proves \eqref{eq:actual-tilt-KL}.
Pinsker's inequality with $\TV=\sup_E|P(E)-Q(E)|$ gives
$\TV^2\le D_{\rm KL}/2$. Finally,
$Q_0\le C_\beta\int f^2d\nu$ by the actual intensity bound.
\end{proof}

\begin{corollary}[The mean measure is controlled as well]
\label{cor:actual-tilt-intensity}
For a relatively compact Borel window $A$, put
$m_0=\mathbb E_{N,\beta}\Xi_N(A)$ and
$D=\frac{|t|}{2}e^{|t|M/2}\sqrt{Q_0}$.
The variation of the two mean measures restricted to $A$ is at most
\[
 2\sqrt{m_0}\,D+D^2.
\]
Here variation is the supremum of the integral difference over
all real $u$ supported in $A$ with $|u|\le1$.
\end{corollary}
\begin{proof}
Write $m(s)=\mathbb E_s\Xi_N(A)$. The susceptibility estimate gives
$|m'(s)|\le\sqrt{m(s)Q(s)}$. Applying the inequality first to
$\sqrt{m(s)+\epsilon}$ and then letting $\epsilon\downarrow0$
shows
$\sqrt{m(s)}\le\sqrt{m_0}+\tfrac12\int_0^{|s|}\sqrt{Q(\pm v)}dv$,
where the sign follows that of $s$.
For $u$ as stated,
$|\frac{d}{ds}\mathbb E_s\Xi_N(u)|\le\sqrt{m(s)Q(s)}$.
Integrate this bound and use
$\tfrac12\int_0^{|t|}\sqrt{Q(\pm v)}dv\le D$.
The result is $2\sqrt{m_0}D+D^2$, independently of $u$.
\end{proof}

\subsection{A distant Green tail may have large mean but small effect}

Let
\[
 g(a,z)=-\log\left|\frac{a-z}{1-a\bar z}\right|.
\]
Fix $0<r<1$, let $\sigma$ be a finite real signed source measure
supported in $\overline B_r$, and write $q=|\sigma|(\overline B_r)$.
For $R>r$ put
\[
 V_R(z)=\mathbf1_{\{|z|>R\}}\int g(a,z)\,\sigma(da),
 \qquad \delta=1-R^2.
\]
This is a potential applied only to the distant particles; $t$ below
is its fixed amplitude, not an $N$-dependent temperature change.

\begin{theorem}[Full-law attenuation of distant Green-potential tilts]
\label{thm:green-tail-TV}
Suppose $0<\beta\le2$ and $\delta\le(1-r)^2/2$. For every $N\ge1$,
every real $t$, and every source just described,
\begin{equation}\label{eq:green-tail-TV}
 \TV(P_{N,\beta}^{t,V_R},P_{N,\beta})
 \le \frac{|t|q\sqrt{C_\beta}}{2(1-r)^2}
 \exp\!\left\{\frac{|t|q\delta}{2(1-r)^2}\right\}\sqrt\delta.
\end{equation}
The same rate controls every local law, including every particle
count. On each fixed compact window it also controls the mean
measure at order $O_{\beta,r,t,q,A}(\sqrt\delta)$.
\end{theorem}
\begin{proof}
The identity
\[
 g(a,z)=-\frac12\log(1-u),\qquad
 u=\frac{(1-|a|^2)(1-|z|^2)}{|1-a\bar z|^2}
\]
and the assumptions give $0\le u\le1/2$. Since
$-\frac12\log(1-u)\le u$ on this interval,
\[
 |V_R(z)|\le\frac{q}{(1-r)^2}(1-|z|^2)\mathbf1_{\{|z|>R\}}.
\]
Consequently
\[
 \|V_R\|_\infty\le\frac{q\delta}{(1-r)^2},\qquad
 \int V_R^2d\nu\le\frac{q^2\delta}{(1-r)^4},
\]
because $(1-|z|^2)^2d\nu=d\mu$ and
$\mu(\{|z|>R\})=\delta$. Apply
Theorem~\ref{thm:actual-tilt-TV}; restriction is a measurable map and
cannot increase TV. The mean-measure assertion follows from
Corollary~\ref{cor:actual-tilt-intensity} and
$m_0\le C_\beta\nu(A)$.
\end{proof}

The estimate does not follow from a small supremum norm alone:
the elementary energy bound would be $N\|V_R\|_\infty$, which is
not uniform in $N$. Moreover $\int |V_R|d\nu$ need not be finite.
The proof removes the deterministic mean through normalization and
controls the fluctuating sum by its quadratic intensity cost.
In curvature $-1$, $\sqrt\delta=\operatorname{sech}
(d(0,R)/2)$; the rate is exponential in the hyperbolic distance
at fixed source region. No boundary-volume factor has been omitted.

\subsection{The entire distant Green force vanishes in mean square}

For $|a|<R$ define the complex force
\[
 \mathcal E_{N,R}(a)
 :=\sum_{|z_i|>R}\frac{1-|z_i|^2}{(z_i-a)(1-a\bar z_i)}
 =2\partial_a\sum_{|z_i|>R}g(a,z_i).
\]
This is the field from all particles beyond $R$, without a truncation
of its angular modes.

\begin{theorem}[Screening and reconstruction of the Green force]
\label{thm:green-force-reconstruction}
For $0<\beta\le2$, $0<r<1$ and $R\ge(1+r)/2$,
\begin{equation}\label{eq:green-force-L2}
 \mathbb E\mathcal E_{N,R}(a)=0,\qquad
 \sup_{|a|\le r}\mathbb E|\mathcal E_{N,R}(a)|^2
 \le\frac{4C_\beta(1-R^2)}{(1-r)^4},
 \qquad N\ge1.
\end{equation}
For every fixed $r<1$ there is also $C_{\beta,r}'<\infty$ such that
\[
 \mathbb E\sup_{|a|\le r}|\mathcal E_{N,R}(a)|^2
 \le C_{\beta,r}'(1-R^2)
\]
for all sufficiently large $R$, uniformly in $N$.

Let $\Xi_{N_j}$ converge locally in law to a subsequential limit
$\Xi$. For any $0<s<1$, the exterior force
\begin{equation}\label{eq:configuration-force}
 \mathcal E_\Xi^{s}(a)
 =\lim_{R\uparrow1}
 \sum_{z\in\Xi:\ s<|z|<R}
       \frac{1-|z|^2}{(z-a)(1-a\bar z)},\qquad |a|<s,
\end{equation}
exists in mean square, locally uniformly in $a$. It is measurable
from the exterior configuration $\Xi|_{\D\setminus B_s}$.
Moreover the finite-gas exterior fields $\mathcal E_{N_j,s}$
converge jointly with the configurations to this field. Thus this
Green force carries no additional random contribution from infinity.
\end{theorem}
\begin{proof}
For a fixed radius $u>|a|$, angular averaging gives
\[
 \frac1{2\pi}\int_0^{2\pi}g(a,ue^{i\theta})d\theta=-\log u.
\]
Indeed the mean of $\log|a-ue^{i\theta}|$ is $\log u$, and that
of $\log|1-au e^{-i\theta}|$ is zero, by their convergent harmonic
series. Differentiation in $a$ and radiality of the finite-gas
intensity prove the zero mean in \eqref{eq:green-force-L2}.

For $|a|\le r$ and $|z|>R\ge(1+r)/2$, the force summand $f_a(z)$
satisfies
\[
 |f_a(z)|\le\frac{2(1-|z|^2)}{(1-r)^2}.
\]
Apply \eqref{eq:subpoisson-variance} to its real and imaginary
parts and add the bounds. The zero mean then gives
\[
 \mathbb E|\mathcal E_{N,R}(a)|^2
 \le\int_{|z|>R}|f_a(z)|^2\rho_{N,\beta}(z)d\mu(z),
\]
which proves \eqref{eq:green-force-L2} by the intensity bound.
Choose $r<r_1<1$. The field is holomorphic on $B_R$. The disk
mean-value inequality for its squared modulus implies
\[
 \sup_{B_r}|\mathcal E_{N,R}|^2
 \le\frac{4}{(r_1-r)^2}\int_{B_{r_1}}
                     |\mathcal E_{N,R}(a)|^2d\mu(a).
\]
For $R\ge(1+r_1)/2$, the already proved estimate and Tonelli give
the locally uniform bound. The identical proof applies to a finite
radial annulus $R<|z|<R'$ and gives the same upper bound with $R$.

We detail passage to the point-process limit. On a compact window,
\eqref{eq:subpoisson-generating} with $v=e^{u\mathbf1_A}$ and
\eqref{eq:actual-uniform-bergman-upper} bound every exponential
count moment uniformly in $N$. Local convergence therefore passes
expectations of continuous compactly supported linear statistics,
and their moments, to the limit. In particular, the limiting mean
measure is rotation invariant and bounded by $C_\beta\nu$;
each fixed circle has zero mean mass and hence almost surely contains no
points. These observations justify restriction to each annulus
$s<|z|<R$ and convergence of its force, locally uniformly in $a$.
One may first work with smooth radial cutoffs and then remove
them using the uniform count-moment and intensity bounds.

The annular mean-square estimate just proved passes to the limit.
For every $r<s$, its locally uniform version makes the sums in
\eqref{eq:configuration-force} Cauchy in
$L^2(\sup_{B_r}|\cdot|)$. Their limit is a measurable function of
the exterior points. Choose a countable sequence $r\uparrow s$;
the limits agree on overlaps and define the asserted holomorphic
field. Finally, truncate both the finite-gas and limiting fields
at radius $R$. Their joint convergence with configurations holds
for fixed $R$, while the discarded tails have uniformly vanishing
mean-square norms as $R\uparrow1$. The triangle inequality for
bounded Lipschitz tests proves the stated joint convergence.
\end{proof}

Integrating this force from zero to $a$ shows that the Green potential
modulo its value at zero is recovered as
\[
 H_\Xi^s(a)=\lim_{R\uparrow1}
       \sum_{z\in\Xi:\ s<|z|<R}[g(a,z)-g(0,z)],\qquad |a|<s,
\]
with the same local mean-square convergence. Every finite annular
sum has mean zero by radiality. This removes the need to specify a
separate random mark for this particular Green force. It does not
identify the additive potential constant, the conditional count
weights, or the harmonic part of the original Euclidean interaction.
Indeed
\[
 -\log|a-z|=g(a,z)-\log|1-a\bar z|,
\]
and the second term has boundary modes that the present estimate
does not suppress. Consequently
Theorem~\ref{thm:green-force-reconstruction} is not a proof of independence
from the complete distant exterior sigma-field.


\section{The full actual exterior at the determinantal benchmark}
\label{sec:full-exterior-benchmark}

The Green-tail tilt estimates do not themselves bound conditioning on
the entire far exterior.  At $\beta=2$, however, that stronger statement
has a direct proof.  It is uniform in the finite particle number, and
therefore handles exterior data chosen by the actual gases even along
triangular arrays of radii and particle numbers.

For random elements $X,Y$ of standard Borel spaces, write
\[
 \mathcal B(X,Y)
 =d_{\mathrm{TV}}\bigl(\operatorname{Law}(X,Y),
              \operatorname{Law}(X)\otimes\operatorname{Law}(Y)\bigr),
\]
where probability total variation takes values in $[0,1]$.
The disintegration formula for total variation gives
\begin{equation}\label{eq:full-exterior-disintegration}
 \mathcal B(X,Y)
 =\mathbb E\,d_{\mathrm{TV}}
       \bigl(\operatorname{Law}(X\mid Y),\operatorname{Law}(X)\bigr)
 =\mathbb E\,d_{\mathrm{TV}}
       \bigl(\operatorname{Law}(Y\mid X),\operatorname{Law}(Y)\bigr).
\end{equation}
For clarity, this identity follows by disintegrating the signed measure
$\operatorname{Law}(X,Y)-\operatorname{Law}(X)\otimes\operatorname{Law}(Y)$
over either marginal.  A countable generating algebra for the standard
Borel space, followed by monotone approximation, makes the fiberwise
variation measurable and identifies the variation of the signed measure
with its integral.  Consequently the order of conditioning in
\eqref{eq:full-exterior-disintegration} can be reversed without changing
the dependence coefficient.

\begin{theorem}[Entire selected-exterior attenuation at two]
\label{thm:full-exterior-beta-two}
Let $\Xi_N$ have density proportional to $|\Delta_N|^2$ relative to
$\mu^N$, where $\mu=dA/\pi$ on the unit disk.  Fix $0<r<R<1$, and put
$X_N=\Xi_N|_{B_r}$ and $Y_{N,R}=\Xi_N|_{\mathbb D\setminus B_R}$.
Then
\begin{align}
 \mathcal B(X_N,Y_{N,R})
 &\le \min\left\{1,
   2\sum_{j=1}^N
       \frac{r^{2j}(1-R^{2j})}{1-r^{2j}}\right\}
                                                        \label{eq:full-exterior-exact}\\
 &\le 2(1-R^2)\sum_{j=1}^{\infty}
                      \frac{j r^{2j}}{1-r^{2j}}
 \le \frac{2r^2(1-R^2)}{(1-r^2)^3}.
                                                        \label{eq:full-exterior-rate}
\end{align}
In particular,
\[
 \lim_{R\uparrow1}\sup_{N\ge1}
 \mathbb E\,d_{\mathrm{TV}}
  \bigl(\operatorname{Law}(\Xi_N|_{B_r}\mid
                    \Xi_N|_{\mathbb D\setminus B_R}),
                    \operatorname{Law}(\Xi_N|_{B_r})\bigr)=0.
\]
This controls the whole local configuration, including its particle
number, and every measurable statistic of the whole far exterior.
\end{theorem}

We record a proof of the finite-projection comparison needed here.
It is a special case of the ordered-kernel theorem
\cite[Theorem 3.8]{LyonsScreening}; its proof below removes any need
for an external coupling theorem in the quantitative argument.

\begin{lemma}[Finite projection comparison]
\label{lem:finite-projection-comparison}
Let $H_1\subset H_2$ be finite-dimensional subspaces of $L^2(E,\mu)$,
and let $Q_i$ be their projection determinantal laws.  Then
$\mathbb E_{Q_1}F\le\mathbb E_{Q_2}F$ for every bounded measurable
function $F$ increasing under inclusion of configurations.  For every
measurable $B\subset E$,
\begin{equation}\label{eq:ordered-projection-local-TV}
 d_{\mathrm{TV}}(Q_1|_B,Q_2|_B)
 \le\mathbb E_{Q_2}\Xi(B)-\mathbb E_{Q_1}\Xi(B).
\end{equation}
\end{lemma}
\begin{proof}
It suffices first to compare ranks $n$ and $n+1$, with orthonormal
basis $\phi_1,\ldots,\phi_n$ of $H_1$ and unit vector
$\phi\perp H_1$ such that $H_2=H_1\oplus\mathbb C\phi$.
In the antisymmetric subspace of $L^2(E^n,\mu^n)$, the normalized
wave function
\[
 \Psi(z_1,\ldots,z_n)=\frac1{\sqrt{n!}}
                      \det[\phi_i(z_j)]_{i,j=1}^n
\]
has squared modulus equal to the ordered density of $Q_1$.
For an antisymmetric function $f$ of $n$ variables define
\[
 (c_\phi f)(z_1,\ldots,z_{n+1})
 =\frac1{\sqrt{n+1}}\sum_{i=1}^{n+1}(-1)^{i-1}\phi(z_i)
                        f(z_1,\ldots,\widehat z_i,\ldots,z_{n+1}),
\]
where the hat means that the coordinate is omitted.  Its adjoint on
antisymmetric functions is the operator
\[
 (a_\phi f)(z_1,\ldots,z_{n-1})
   =\sqrt n\int_E\overline{\phi(t)}f(t,z_1,\ldots,z_{n-1})\,d\mu(t).
\]
Expanding the finite sum and using antisymmetry gives
\[
 a_\phi c_\phi+c_\phi a_\phi=I,
 \qquad
 \|c_\phi f\|_2^2=\|f\|_2^2-\|a_\phi f\|_2^2\le\|f\|_2^2.
\]
For example, in the first identity the term inserting $t$ in the first
coordinate contributes $f$, while the remaining $n$ terms are exactly
$-c_\phi a_\phi f$ after their permutation signs are collected.  The
norm identity follows by taking the inner product with $f$ and using
adjointness.  These calculations also apply to $n=0$ with the scalar
vacuum space and zero annihilation operator.

Orthogonality implies $a_\phi\Psi=0$.  The normalized wave function
of $Q_2$ is $\Psi'=c_\phi\Psi$, up to an irrelevant sign, and
$a_\phi\Psi'=\Psi$.  Let $\mathcal U$ be an increasing measurable
configuration event, and let $A_n,A_{n+1}$ be its symmetric ordered
sectors of cardinalities $n,n+1$.  Put $f=\mathbf1_{A_n}\Psi$.
The function $c_\phi f$ vanishes off $A_{n+1}$: any nonzero summand
exhibits a subconfiguration belonging to $\mathcal U$.  Hence, with
$p_i=Q_i(\mathcal U)$,
\[
 p_1=\langle\Psi,f\rangle
     =\langle\Psi',c_\phi f\rangle
     =\langle\mathbf1_{A_{n+1}}\Psi',c_\phi f\rangle
     \le\sqrt{p_2}\,\|c_\phi f\|_2
     \le\sqrt{p_2p_1}.
\]
The inner product is real and nonnegative because its value is $p_1$;
Cauchy--Schwarz justifies the displayed inequality.  Thus $p_1\le p_2$.
Add an orthonormal basis of $H_2\ominus H_1$ one vector at a time and
then use the layer-cake formula to obtain the assertion for every
bounded increasing measurable function.

For a measurable event $F$ of the restriction to $B$, both functions
\[
 \eta\longmapsto\eta(B)+\mathbf1_F(\eta|_B),
 \qquad
 \eta\longmapsto\eta(B)-\mathbf1_F(\eta|_B)
\]
are increasing under inclusion.  Indeed, an added point outside $B$
changes neither term; a strict inclusion on $B$ increases its count
by at least one and changes the indicator by at most one.  Counts are
bounded by the finite ranks on the two laws; the increasing functions
can therefore be truncated outside that range when applying the
preceding assertion.  Applying the assertion to both signs gives
\[
 |(Q_2|_B)(F)-(Q_1|_B)(F)|
 \le\mathbb E_{Q_2}\Xi(B)-\mathbb E_{Q_1}\Xi(B).
\]
Taking a supremum over $F$ proves the total-variation inequality.
\end{proof}

\begin{proof}[Proof of Theorem~\ref{thm:full-exterior-beta-two}]
Put $A=B_r$, $E=\mathbb D\setminus A$, and $B=\mathbb D\setminus B_R$.
There are almost surely no particles on either bounding circle.  The
event $\{\Xi_N(A)=0\}$ has positive probability, since the joint density
is positive almost everywhere on $E^N$.  Let $Q_0$ be the conditional
law of the entire configuration in $E$ given that event.

The Vandermonde determinant and Gram orthonormalization identify $Q_0$
as the projection determinantal process onto
\[
 H_0=\operatorname{Pol}_{N-1}
       \quad\hbox{in }L^2(E,\mu).
\]
More generally, condition on the complete configuration
$x=\{x_1,\ldots,x_m\}$ in $A$.  Almost surely its points are distinct,
and $0\le m\le N$.  Factoring the joint Vandermonde shows that the
conditional density of the $N-m$ exterior points is proportional to
\begin{equation}\label{eq:conditioned-exterior-polynomials}
 |\Delta_{N-m}(z_1,\ldots,z_{N-m})|^2
   \prod_{i=1}^{N-m}|p_x(z_i)|^2\prod_{i=1}^{N-m}d\mu(z_i),
 \qquad p_x(z)=\prod_{a\in x}(z-a),\quad z_i\in E.
\end{equation}
Terms depending only on $x$ cancel in the conditional normalization.
The density in \eqref{eq:conditioned-exterior-polynomials} is the
projection determinantal law onto
\[
 H_x=p_x\operatorname{Pol}_{N-m-1}\subseteq H_0
             \quad\hbox{in }L^2(E,\mu).
\]
For $m=N$ the subspace is zero and the exterior is empty.  Otherwise
the identity follows by taking a determinant of any basis of $H_x$;
the factors $p_x$ give exactly the displayed one-body weight.
Distinct interpolation nodes impose $m$ independent conditions on
$\operatorname{Pol}_{N-1}$, so the displayed subspace also consists of
all polynomials in $H_0$ vanishing at every point of $x$.

Write $Q_x$ for this conditional exterior law and $P_x$ for its
projection kernel in $L^2(E,\mu)$.  The subspace inclusion gives
$P_x\le P_0$.  Applying Lemma~\ref{lem:finite-projection-comparison} on $B$ yields
\begin{equation}\label{eq:exterior-hole-bracket}
 d_{\mathrm{TV}}(Q_x|_B,Q_0|_B)
 \le \mathbb E_{Q_0}\Xi(B)-\mathbb E_{Q_x}\Xi(B).
\end{equation}
No regularity or typicality requirement on the observed points $x$ is
needed for this inequality, beyond distinctness and admissibility.

Let $\overline Q=\operatorname{Law}(\Xi_N|_E)$, the mixture of $Q_x$
against the actual law of $X_N$.  Convexity of total variation and
\eqref{eq:exterior-hole-bracket} imply
\[
 d_{\mathrm{TV}}(\overline Q|_B,Q_0|_B)
 \le\mathbb E\,d_{\mathrm{TV}}(Q_{X_N}|_B,Q_0|_B)
 \le \Delta_N(r,R),
\]
where
\[
 \Delta_N(r,R)=
       \mathbb E[\Xi_N(B)\mid\Xi_N(A)=0]-\mathbb E\Xi_N(B).
\]
A second triangle inequality and
\eqref{eq:full-exterior-disintegration} give
\begin{equation}\label{eq:full-exterior-hole-reduction}
 \mathcal B(X_N,Y_{N,R})\le2\Delta_N(r,R).
\end{equation}

It remains to compute this scalar excess count.  Angular orthogonality
in $E$ gives the orthonormal basis
\[
 \left\{\sqrt{\frac{j}{1-r^{2j}}}\,z^{j-1}:1\le j\le N\right\}
\]
of $H_0$.  Thus
\[
 \mathbb E[\Xi_N(B)\mid\Xi_N(A)=0]
   =\sum_{j=1}^N\frac{1-R^{2j}}{1-r^{2j}},
 \qquad
 \mathbb E\Xi_N(B)=\sum_{j=1}^N(1-R^{2j}).
\]
Subtracting proves the sum in \eqref{eq:full-exterior-exact}.
Finally $1-R^{2j}\le j(1-R^2)$ and
$1-r^{2j}\ge1-r^2$.  Summing
$\sum_{j\ge1}j r^{2j}=r^2/(1-r^2)^2$ proves
\eqref{eq:full-exterior-rate}.
\end{proof}

\begin{corollary}[Typical exteriors and normalized boundary changes]
\label{cor:full-exterior-typical}
Let $b_r(R)$ denote the last expression in
\eqref{eq:full-exterior-rate}.  For every $\varepsilon>0$,
\[
 \mathbb P\!\left\{
 d_{\mathrm{TV}}\bigl(\operatorname{Law}(X_N\mid Y_{N,R}),
                     \operatorname{Law}(X_N)\bigr)>\varepsilon\right\}
 \le \frac{b_r(R)}{\varepsilon}.
\]
Suppose a changed probability law has density $W(Y_{N,R})$ relative
to the actual finite gas, where $W\ge0$ and $\mathbb EW=1$.  If
$1<p\le\infty$ and $q=p/(p-1)$, then
\begin{equation}\label{eq:boundary-likelihood-local-TV}
 d_{\mathrm{TV}}\bigl(\operatorname{Law}_{W}(X_N),
                     \operatorname{Law}(X_N)\bigr)
 \le \|W\|_{L^p}\,b_r(R)^{1/q}.
\end{equation}
For $p=\infty$, the convention is $q=1$.  Thus bounded normalized
likelihoods, and more generally uniformly $L^p$ normalized likelihoods,
give a rigorously admissible class of boundary changes.
\end{corollary}
\begin{proof}
The first assertion is Markov's inequality.  Set
$d(Y)=d_{\mathrm{TV}}(\operatorname{Law}(X_N\mid Y),
\operatorname{Law}(X_N))$.  For each local event $F$, conditioning on
$Y$ gives
\[
 |\mathbb E[W\mathbf1_F]-\mathbb P(F)|
 \le\mathbb E[W d(Y)].
\]
Take a supremum over $F$, use H\"older's inequality, and use
$0\le d\le1$ and $\mathbb Ed\le b_r(R)$.
\end{proof}

\begin{remark}[What the benchmark excludes and what it does not]
The estimate is uniform in $N$ with $r$ fixed.  Consequently no
measurable estimator built from the actual far exterior can preserve a
positive amount of total-variation information about the whole fixed
local configuration along $R\uparrow1$, even if $N$ changes arbitrarily.
This includes exterior harmonic sums, exterior counts, and functions
whose formula depends on $N$ and $R$.  It is stronger than invoking
tail triviality of the limiting Bergman process after taking $N$ to
infinity first.

The constant is not uniform as the observation radius $r$ tends to one.
Nor is the conditional distance small for every individual exterior
configuration: the theorem gives an average and hence a high-probability
bound under the selected boundary law.  A rare boundary event can have
a normalized likelihood with an arbitrarily large $L^p$ norm in
\eqref{eq:boundary-likelihood-local-TV}.  Therefore the result is
consistent with adverse deterministic exterior configurations.
\end{remark}

\subsection{A formulation transferable beyond determinants}

The reversal of conditioning in the preceding proof is a probability
argument, not a determinantal one.  It gives a concrete alternative
target for screening at other temperatures.

\begin{proposition}[Full-law attenuation from a hole response]
\label{prop:general-hole-response}
Let $\Xi$ be any finite point process on a standard Borel region $D$,
and let $A,B\subset D$ be disjoint measurable sets.  Assume
$\mathbb P\{\Xi(A)=0\}>0$ and that all the expectations below are finite.
Write $Q_0$ for the law of $\Xi|_{D\setminus A}$ conditioned on the
hole event and $Q_x$ for its law conditioned on the complete
configuration $\Xi|_A=x$.  Suppose that, for almost every $x$ under the
actual interior marginal, $Q_x|_B$ is stochastically dominated by
$Q_0|_B$ under inclusion; an inclusion coupling is a sufficient
condition.
Then
\begin{equation}\label{eq:general-hole-response}
 \mathcal B(\Xi|_A,\Xi|_B)
 \le2\left\{
  \mathbb E[\Xi(B)\mid\Xi(A)=0]-\mathbb E\Xi(B)\right\}.
\end{equation}
More generally, it suffices that for some fixed reference law $Q_*$
and measurable $e(x)\ge0$,
\[
 d_{\mathrm{TV}}(Q_x|_B,Q_*|_B)\le e(x),
 \qquad \mathbb E e(\Xi|_A)<\infty.
\]
Then $\mathcal B(\Xi|_A,\Xi|_B)\le2\mathbb E e(\Xi|_A)$.
\end{proposition}
\begin{proof}
Apply stochastic domination to the two increasing functions
$\eta(B)\pm\mathbf1_F(\eta|_B)$, truncated if necessary and
then passed to the limit by the assumed finite count expectations.
As in the proof of Lemma~\ref{lem:finite-projection-comparison}, this
bounds the distance to $Q_0|_B$ by its excess expected count.  Integrate over the actual
interior configuration.  Convexity bounds the distance between the
unconditional exterior and the same reference law by that integral.
The triangle inequality and
\eqref{eq:full-exterior-disintegration} prove both assertions.
\end{proof}

For the actual disk gas this proposes two separate non-determinantal
estimates: an inclusion coupling, or a controlled approximate
replacement, for complete interior conditioning; and vanishing
far-exterior excess count after imposing a fixed inner hole.  The
susceptibility inequality at $0<\beta\le2$ alone does not establish
either assertion.  In particular, a two-particle-to-one-particle
deletion coupling cannot be iterated over arbitrary particle numbers
without a further theorem.  Proposition~\ref{prop:general-hole-response}
is thus a proved reduction, not a claim of full screening below two.

% Bibliography entry required by the parent document:
% \bibitem{LyonsScreening}
% R. Lyons, Determinantal probability: basic properties and conjectures,
% Proceedings of the International Congress of Mathematicians 2014,
% vol. IV, 137--161. \href{https://arxiv.org/abs/1406.2707}{arXiv:1406.2707}.
% The ordered-kernel theorem used here is Theorem 3.8.


\section{Why susceptibility alone does not give deletion couplings}
\label{sec:lorentzian-covering-obstruction}

The susceptibility estimate does not, by itself, justify replacing an
exterior conditional law by a configuration containing all its particles.
The following example isolates the logical gap.  It is an example of a
discrete probability measure, not a counterexample for the disk gas.

\begin{proposition}[A Lorentzian measure without deletion domination]
\label{prop:lorentzian-no-covering}
There is a strictly positive probability measure on every three-element
subset of a five-element set whose generating polynomial is Lorentzian,
but which fails stochastic covering after conditioning on one coordinate.
\end{proposition}

\begin{proof}
Use coordinates $x_1,x_2,x_3,y,z$, and put
\[
 F=e_3(x_1,x_2,x_3,y,z)+3x_1x_2x_3,
\]
where $e_3$ is the sum of all squarefree monomials of degree three.
Thus the coefficient of $x_1x_2x_3$ is four and each of the other nine
coefficients is one.  The support is the set of all three-element
subsets, which has the matroid exchange property.

For a homogeneous cubic with this support, the Lorentzian criterion
used above asks that each first derivative have a Hessian with at most
one positive eigenvalue.  The derivatives with respect to $y$ and $z$
are $e_2$ on the four remaining variables.  Their nonzero Hessian block
is the matrix with zero diagonal and all off-diagonal entries one,
whose eigenvalues are $3,-1,-1,-1$.

For the derivative with respect to $x_1$, order the remaining variables
as $x_2,x_3,y,z$.  Its Hessian block is
\[
 H=\begin{pmatrix}
 0&4&1&1\\
 4&0&1&1\\
 1&1&0&1\\
 1&1&1&0
 \end{pmatrix}.
\]
The vectors $(1,-1,0,0)$ and $(0,0,1,-1)$ have eigenvalues $-4$ and
$-1$.  On vectors $(a,a,b,b)$ the action is represented by
$\left(\begin{smallmatrix}4&2\\2&1\end{smallmatrix}\right)$,
whose eigenvalues are $5$ and $0$.  The derivatives with respect to
$x_2$ and $x_3$ have the same spectrum.  The unused differentiated
coordinate adds only a zero eigenvalue.  This verifies the complete
Lorentzian criterion.

Normalize the coefficients of $F$ to a probability measure $\mu$.
Since $F(1,1,1,1,1)=13$, and writing $I_y,I_z$ for occupancy indicators,
\[
 \mu(I_y=1)=\mu(I_z=1)=\frac6{13},\qquad
 \mu(I_y=I_z=1)=\frac3{13}.
\]
Consequently
\[
 \operatorname{Cov}_{\mu}(I_y,I_z)
 =\frac3{13}-\frac{36}{169}=\frac3{169}>0.
\]
More explicitly, the law on the other four coordinates given $I_y=1$
has two particles and assigns probability $1/2$ to $I_z=1$; the law
given $I_y=0$ has three particles and assigns probability $3/7$ to that
same increasing event.  The two-particle configuration therefore cannot
be a subset of the three-particle configuration under any coupling.
\end{proof}

\begin{remark}
This example is not asserted to have coefficients of the form
$\prod_{\{i,j\}\subset S}|a_i-a_j|^\beta$ for planar points $a_i$.
It disproves the implication from the already established Lorentzian
property to deletion domination.  Additional structure of the disk
interaction could still imply the latter.
\end{remark}

For precision, the additional property has a finite system of exact
inequalities. Fix $\beta>0$ and $1\le n\le q$.
Let $a_1,\ldots,a_q$ be distinct points, let $y$ be a
further point, and let $w_i>0$.  For a subset $T\subset\{1,\ldots,q\}$
of size $n$ and a subset $S$ of size $n-1$, set
\[
 A_T=\prod_{i\in T}w_i
       \prod_{\{i,j\}\subset T}|a_i-a_j|^\beta,
 \qquad
 B_S=\prod_{i\in S}\bigl(w_i|a_i-y|^\beta\bigr)
       \prod_{\{i,j\}\subset S}|a_i-a_j|^\beta.
\]
Write $Z_A=\sum_{|T|=n}A_T$ and $Z_B=\sum_{|S|=n-1}B_S$.
For a family $\mathcal C$ of $(n-1)$-element subsets, define its upper
shadow by
\[
 \partial^+\mathcal C
 =\{T:|T|=n,\ S\subset T\text{ for some }S\in\mathcal C\}.
\]

\begin{proposition}[The exact finite deletion criterion]
\label{prop:finite-deletion-hall}
A coupling of $Q(S)=B_S/Z_B$ and $P(T)=A_T/Z_A$ supported on $S\subset T$
exists if and only if, for every family $\mathcal C$,
\begin{equation}
 Z_A\sum_{S\in\mathcal C}B_S
 \le Z_B\sum_{T\in\partial^+\mathcal C}A_T.
 \label{eq:finite-deletion-hall}
\end{equation}
\end{proposition}

\begin{proof}
Necessity follows because $S\in\mathcal C$ and $S\subset T$ imply
$T\in\partial^+\mathcal C$.  For sufficiency, build a finite network
with a source, one vertex for each $S$, one vertex for each $T$, and a
sink.  Give source-to-$S$ edges capacity $Q(S)$, $T$-to-sink edges
capacity $P(T)$, and each inclusion edge $S\to T$ capacity two.
A cut of capacity less than one cannot cross an inclusion edge.
If its source side contains a family $\mathcal C$ of lower vertices,
it must therefore contain their upper shadow, and its capacity is at
least $1-Q(\mathcal C)+P(\partial^+\mathcal C)\ge1$.
All other cuts also have capacity at least one.  The max-flow/min-cut
theorem gives a flow of value one.  Since both sets of marginal
capacities sum to one, all source and sink edges are saturated, and
the flows on inclusion edges form the required coupling.
\end{proof}

For $n=2$, the families $\mathcal C$ are just sets of possible singleton
locations; the Hilbert-space calculation in the two-to-one deletion
theorem proves all these inequalities.  For $n\ge3$, families of
multi-particle configurations must be considered.  Linear-statistic
variance bounds do not check those families.  Establishing
\eqref{eq:finite-deletion-hall} for all relevant planar distance weights,
or a suitably restricted version for boundary data selected by the
actual gas, remains an explicit possible route to a count-sensitive
comparison.


\section{An exact finite-range splitting and a uniformly controlled reference gas}
\label{sec:exact-smearing-reference}

This section concerns the actual unweighted disk gas.  It supplies an exact
representation, including every particle-number weight.  The representation
does not by itself prove screening: its final, positive quadratic tilt still
has to be controlled.  All measures and kernels used below are constructed
explicitly.

Put $\mu=\pi^{-1}dA$, $\nu=\pi^{-1}(1-|z|^2)^{-2}dA$, and let $d$ denote
hyperbolic distance for curvature $-1$.  Write
\[
 \begin{gathered}
 C(z,w)=-\log|z-w|=g(d(z,w))+h(z,w),\\
 g(r)=-\log\tanh(r/2),\qquad h(z,w)=-\log|1-z\bar w|.
 \end{gathered}
\]
The normalized radial element of $\nu$ about any point is
$\frac12\sinh r\,dr$.  For $a>0$, let $S_a$ average a function over the
hyperbolic sphere of radius $a$, with its rotational probability measure.
This operator preserves harmonic functions.  It preserves $\nu$ and is
self-adjoint on $L^2(\nu)$: the invariant measure on pairs of points at
distance $a$ is unchanged on exchanging the points.  Source and target
averaging of an invariant radial kernel commute.

\begin{proposition}[Exact finite-range decomposition]
\label{prop:smearing-exact}
Let $C_a$ be $C$ averaged by $S_a$ in each argument.  There is a nonnegative
radial kernel $k_a$ such that
\[
 C(z,w)=C_a(z,w)+k_a(d(z,w)),\qquad
 k_a(r)=0\quad(r\ge2a).
\]
More precisely, with $q_a(r)=(g(r)-g(a))_+$,
\begin{align}
 k_a&=q_a+S_aq_a,\\
 M_a:=\int k_a(d(o,z))\,\nu(dz)&=2\log\cosh(a/2),\\
 C_a(z,z)&=g(a)-\log(1-|z|^2).
\end{align}
The kernel $C_a$ is finite and positive semidefinite.  For all
$\beta>0$ and $\alpha\ge0$,
\begin{equation}\label{eq:short-bond-norm}
 \int e^{\alpha d(o,z)}\bigl(1-e^{-\beta k_a(d(o,z))}\bigr)\nu(dz)
 \le \beta e^{2\alpha a}M_a.
\end{equation}
\end{proposition}
\begin{proof}
The spherical mean of the Green potential is
\[
 S_ag(r)=g(\max\{r,a\}).
\]
Outside the sphere this is the harmonic mean-value property.  Inside,
the sphere potential is a rotationally invariant harmonic function regular
at its center, hence constant; its value at the center is $g(a)$.
These assertions agree continuously on the sphere, where the logarithmic
singularity in the angular integral is integrable.  Since $h$ is harmonic
in each argument, its spherical means equal $h$.  Therefore
$C_a=S_a^2g+h$ and
$g-S_a^2g=(g-S_ag)+S_a(g-S_ag)=q_a+S_aq_a$.
This proves positivity and the range statement by the triangle inequality.
Using $g'(r)=-1/\sinh r$ and integration by parts gives
\[
 \int q_a\,d\nu
 =\frac12\int_0^a(g(r)-g(a))\sinh r\,dr
 =\frac12\int_0^a\frac{\cosh r-1}{\sinh r}\,dr
 =\log\cosh(a/2).
\]
Measure preservation by $S_a$ gives the stated value of $M_a$.
Furthermore
\[
 S_a^2g(0)=S_a[g(\max\{r,a\})](0)=g(a),
\]
which proves the diagonal formula.

For completeness, positivity as a quadratic form is a different assertion
from pointwise positivity of $k_a$.  For any real finite signed measure
$\sigma$ of finite Green energy, $\iint g\,d\sigma\,d\sigma\ge0$.
One way to verify this is to first use a smooth charge density compactly
supported in a disk $B_R$: its zero-boundary Green potential $U$ satisfies
$-\Delta U=2\pi\sigma$, and integration by parts gives the energy
$(2\pi)^{-1}\int_{B_R}|\nabla U|^2dA\ge0$.  Increasing $R$ to $1$
gives the full-disk assertion; approximating the charges gives it for
finite-energy measures.  In the present application the charges are
finite linear combinations of uniform circle measures, whose logarithmic
energies are finite.  Their energies converge under ordinary smooth
approximation, because the only local singularity is logarithmic.
Thus $S_a^2g$ is positive semidefinite on point configurations.  Also
\[
 \sum_{i,j}t_it_jh(z_i,z_j)
 =\sum_{k\ge1}\frac1k\left|\sum_i t_i z_i^k\right|^2\ge0
 \qquad(t_i\in\mathbb R).
\]
The series converges for every finite configuration in $\mathbb D$.
This proves the assertion for $C_a$.  Finally,
$1-e^{-\beta k_a}\le\beta k_a$ and its support is contained in
$B_{2a}$, proving \eqref{eq:short-bond-norm}.
\end{proof}

For $\eta=\sum_{i=1}^N\delta_{z_i}$ at distinct points, the proposition
gives the following identity without any approximation:
\begin{equation}\label{eq:actual-smearing-density}
 |\Delta_N(\mathbf z)|^\beta\,\mu^N(d\mathbf z)
 =e^{\beta Ng(a)/2}
  e^{-\frac\beta2 C_a(\eta,\eta)}
  e^{-\beta\sum_{i<j}k_a(d(z_i,z_j))}
  \prod_i(1-|z_i|^2)^{2-\beta/2}\nu(dz_i).
\end{equation}
Here $C_a(\sigma,\tau)=\iint C_a(z,w)\,\sigma(dz)\tau(dw)$ whenever
the integral is defined.  The one-body factor is the retained diagonal
self-energy, together with the change of area measure.  It cannot be
dropped when comparing particle counts.

\subsection{Construction of the reference density}

We now assume $0<\beta<2$, and set
\[
 p=2-\beta/2>1,\qquad c=\frac{2p}{\beta}=\frac4\beta-1>0,
 \qquad u(z)=(1-|z|^2)^p.
\]
The appearance of $c$ below concerns a reference density, not the
unproved intensity of the actual particle process.

\begin{proposition}[A finite-mass reference at every activity]
\label{prop:saddle-reference}
For every $\lambda>0$ there exists a positive radial smooth function
$b_\lambda$ on $\mathbb D$ satisfying
\begin{equation}\label{eq:reference-saddle}
 b_\lambda(z)=\lambda u(z)
       \exp[-\beta(C*b_\lambda)(z)],\qquad
 (C*b)(z)=\int C(z,w)b(w)\nu(dw).
\end{equation}
It has finite mass, is radially decreasing, and satisfies
\[
 0<b_\lambda(z)<c,\qquad
 b_\lambda(z)\sim\lambda(1-|z|^2)^p\quad(|z|\uparrow1).
\]
These references increase strictly with $\lambda$ and
\begin{equation}\label{eq:saddle-local-limit}
 b_\lambda\longrightarrow c
 \quad\hbox{locally uniformly as }\lambda\longrightarrow\infty.
\end{equation}
Their defining local equation is
\begin{equation}\label{eq:saddle-liouville}
 \Delta_{\mathbb H}\log b_\lambda
       =\frac\beta2(b_\lambda-c).
\end{equation}
\end{proposition}
\begin{proof}
We construct the functions by a radial ordinary differential equation,
which also proves the required bounds without an assumption about the
gas.  Let $t=d(0,z)$, and fix $x\in(0,c)$.  Solve
\begin{equation}\label{eq:saddle-ode}
 v''(t)+\coth t\,v'(t)=\frac\beta2(e^{v(t)}-c),\qquad
 v(0)=\log x,\quad v'(0)=0.
\end{equation}
The regular radial solution exists uniquely near zero.  For example,
the integrated equation
\[
 v(t)=\log x+\frac\beta2\int_0^t\frac1{\sinh s}
                  \int_0^s(e^{v(r)}-c)\sinh r\,dr\,ds
\]
has a unique continuous solution on a short interval by the contraction
mapping theorem, since its double-integral operator has norm $O(t^2)$
on bounded sets.  It then gives the usual regularity and the differential
equation.  As long as the solution exists, the same formula shows that
$v'(t)<0$ for $t>0$ and $e^{v(t)}\le x<c$: a first failure of either
property is impossible.  Moreover
\[
 -p\tanh(t/2)\le v'(t)
 \le-\frac\beta2(c-x)\tanh(t/2).
\]
These bounds prevent finite-time blowup and extend the solution to all
$t\ge0$.  If $b_x(t)=e^{v(t)}$ and
$\delta_x=\beta(c-x)/2>0$, they imply
\begin{equation}\label{eq:saddle-decay-initial}
 x\cosh(t/2)^{-2p}\le b_x(t)
 \le x\cosh(t/2)^{-2\delta_x}.
\end{equation}

An exact refinement of the integrated equation is
\[
 v'(t)=-p\tanh(t/2)
      +\frac\beta{2\sinh t}\int_0^t b_x(s)\sinh s\,ds.
\]
The integral of its second term over $t\in(0,\infty)$ is finite.
Indeed, Tonelli identifies it with
\[
 \frac\beta2\int_0^\infty b_x(s)\sinh s
                      \left(\int_s^\infty\frac{dt}{\sinh t}\right)ds
 =\frac\beta2\int_0^\infty b_x(s)\sinh s\,g(s)\,ds<\infty.
\]
Near zero the integrand is $O(s|\log s|)$; at infinity,
$\sinh s\,g(s)$ is bounded and \eqref{eq:saddle-decay-initial} gives
exponential decay.  Consequently
\begin{align}
 b_x(t)&\sim\lambda(x)\cosh(t/2)^{-2p},\label{eq:saddle-asymptotic}\\
 \lambda(x)&=x\exp\left[\frac\beta2\int_0^\infty
                    b_x(s)\sinh s\,g(s)\,ds\right].
\end{align}
Since $p>1$, \eqref{eq:saddle-asymptotic} proves that $b_x\nu$ has
finite total mass.

Let $B_x=C*b_x$.  Angular averaging gives
$C(z,w)\mapsto-\log\max\{|z|,|w|\}$ for radial charges.  In
particular $B_x\ge0$, it is bounded by $B_x(0)<\infty$, and
$B_x(z)\to0$ as $|z|\uparrow1$.  The latter also follows from
$B_x(r)\le-\log r\int b_x\,d\nu$ for $r>0$.
Its Laplacian is
$\Delta_{\mathbb H}B_x=-b_x/2$, since
$\Delta C(\cdot,w)=-2\pi\delta_w$ in Euclidean coordinates.
Also $\Delta_{\mathbb H}\log(1-|z|^2)=-1$.
It follows from \eqref{eq:saddle-ode} that
$\log b_x-p\log(1-|z|^2)+\beta B_x$ is a radial harmonic
function regular at zero and hence constant.  Taking its boundary
limit using \eqref{eq:saddle-asymptotic} shows that this constant is
$\log\lambda(x)$.  Thus \eqref{eq:reference-saddle} holds for $b_x$
at activity $\lambda(x)$.

It remains to parameterize all activities and prove the limiting claim.
If $0<x<y<c$, the difference $w=v_y-v_x$ starts positive with zero
derivative.  While it is positive,
$(\sinh t\,w')'=\frac\beta2\sinh t(b_y-b_x)>0$ for $t>0$.
It therefore remains positive and increasing.  Thus $b_y>b_x$ and,
by the formula for $\lambda$, $\lambda(y)>\lambda(x)$.
Continuity of $\lambda(x)$ inside $(0,c)$ follows from continuous
dependence of the ODE on initial values and dominated convergence in
the integral formula, using the upper bound in
\eqref{eq:saddle-decay-initial} uniformly on compact $x$-intervals.
As $x\downarrow0$, the same bound, with for example
$\delta_x\ge p/2$, shows that the integral in $\log(\lambda(x)/x)$
tends to zero; hence $\lambda(x)\to0$.
As $x\uparrow c$, the ODE solutions tend on compact intervals to
the constant solution $b=c$.  Monotonicity and Tonelli give
\[
 \int_0^\infty b_x(s)\sinh s\,g(s)\,ds\longrightarrow\infty,
\]
because $\sinh s\,g(s)\to1$.  Therefore $\lambda(x)\to\infty$.
The continuous strictly increasing map $x\mapsto\lambda(x)$ maps
$(0,c)$ onto $(0,\infty)$.  Its inverse defines $b_\lambda$ and
proves all the claims, including local uniform convergence by the
continuous dependence already used.
\end{proof}

\subsection{Exact grandcanonical representation and the small reference norm}

Let $P_{\zeta,\beta}$ denote the grandcanonical mixture of the original
unweighted gases, with unnormalized $N$-particle sector
\[
 \frac{\zeta^N}{N!}|\Delta_N(\mathbf z)|^\beta\prod_i\mu(dz_i).
\]
For $0<\beta<2$ its partition function is finite.  Indeed,
$\int|\Delta_N|^2d\mu^N=N!\prod_{j=0}^{N-1}(j+1)^{-1}=1$
by expanding the determinant and using angular orthogonality.  Jensen's
inequality gives $Z_N(\beta)\le1$, and hence
$\sum_N\zeta^NZ_N(\beta)/N!\le e^\zeta$.

\begin{theorem}[The actual gas as a tilt of a uniformly subcritical reference]
\label{thm:actual-reference-tilt}
Fix $a>0$, $\zeta>0$, and $0<\beta<2$.  Set
\[
 \lambda=\zeta e^{\beta g(a)/2},\qquad b=b_\lambda,\qquad
 a_\lambda(z)=b(z)\exp[\beta(k_a*b)(z)].
\]
Let $Q_\lambda$ be the finite grandcanonical gas with activity
$a_\lambda(z)\nu(dz)$ and pair interaction $\beta k_a$.  Then
\begin{equation}\label{eq:actual-exact-tilt}
 \frac{dP_{\zeta,\beta}}{dQ_\lambda}(\eta)
 =\frac{\exp[-\frac\beta2 C_a(\eta-b\nu,\eta-b\nu)]}
 {E_{Q_\lambda}\exp[-\frac\beta2 C_a(\eta-b\nu,\eta-b\nu)]}.
\end{equation}
The expectation in the denominator is strictly positive.  Uniformly in
$\lambda$,
\begin{align}
 a_\lambda(z)&\le c\exp(\beta cM_a),\label{eq:reference-activity-bound}\\
 \sup_z\int e^{\alpha d(z,w)}a_\lambda(w)
       (1-e^{-\beta k_a(d(z,w))})\nu(dw)
 &\le\beta cM_a\exp(\beta cM_a+2\alpha a).
 \label{eq:reference-influence-bound}
\end{align}
In particular, if $0<\beta\le1/4$ and $a=\sqrt\beta$, the
right-hand side for $\alpha=1$ is at most
$\frac14e^{5/4}<1$.
\end{theorem}
\begin{proof}
The convolution $k_a*b$ is nonnegative and at most $cM_a$, so
$a_\lambda\le e^{\beta cM_a}b$.  Its integral is finite by
Proposition~\ref{prop:saddle-reference}.  Since the pair interaction is
nonnegative, the grand partition function of $Q_\lambda$ lies between
$1$ and $\exp(\int a_\lambda\,d\nu)$, and $Q_\lambda$ has a finite
number of particles almost surely.

All energies in \eqref{eq:actual-exact-tilt} are finite.
Indeed $C*b$ is bounded and $b\nu$ has finite mass.  Radial symmetry
annihilates its nonconstant harmonic moments, so $h*b=0$ and
$C_a*b=S_a^2g*b$, which lies between $0$ and $C*b$ because
$0\le S_a^2g\le g$.  This controls the deterministic and cross energies;
the finite atomic self-energies are covered by
Proposition~\ref{prop:smearing-exact}.  Moreover
$\int[-\log(1-|z|^2)]b(z)\nu(dz)<\infty$ by the boundary asymptotic
$b\sim\lambda(1-|z|^2)^p$ with $p>1$.  Together with
$|h(z,w)|\le\frac12(h(z,z)+h(w,w))$, which follows from the positive
semidefiniteness of $h$ and the arithmetic--geometric mean inequality,
this justifies absolute integration of its harmonic term.  Truncating
the radial background and then passing to the limit consequently
extends the positive-energy formula to this noncompact finite measure.
The full energy of the signed
charge is nonnegative by the same Green-energy argument and harmonic
moment expansion.  Hence the tilt lies in $(0,1]$ on every finite
configuration and has positive expectation.

To prove the identity, apply \eqref{eq:actual-smearing-density} to
every $N$-sector and use
$\log(\lambda u)=\log b+\beta C*b$ and $C=C_a+k_a$.
Completing the square gives the unnormalized equality
\begin{align*}
 &e^{-\frac\beta2 C_a(\eta,\eta)}
   e^{-\beta\sum_{i<j}k_a(d(z_i,z_j))}
   \prod_i\lambda u(z_i)\nu(dz_i)\\
 &\quad=e^{\frac\beta2 C_a(b\nu,b\nu)}
   e^{-\frac\beta2 C_a(\eta-b\nu,\eta-b\nu)}
   e^{-\beta\sum_{i<j}k_a(d(z_i,z_j))}
   \prod_i a_\lambda(z_i)\nu(dz_i).
\end{align*}
This also holds for $N=0$; the two deterministic energy factors then
cancel.  Division by $N!$ and summation over $N$ prove
\eqref{eq:actual-exact-tilt}, including its particle-count weights.

The first bound follows from $b<c$; the second follows from
\eqref{eq:short-bond-norm}.  Finally, $\log\cosh x\le x^2/2$ gives
\[
 M_a=2\log\cosh(a/2)\le a^2/4.
\]
For $a=\sqrt\beta$,
$\beta cM_a\le(4-\beta)\beta/4\le\beta$, so the right-hand side
of \eqref{eq:reference-influence-bound} is at most
$\beta e^{\beta+2\sqrt\beta}\le\frac14e^{5/4}<1$.
\end{proof}

\begin{corollary}[A canonical reference at every particle number]
\label{cor:canonical-reference}
For every $N\ge1$ there is a unique $\lambda_N>0$ such that
$\int b_{\lambda_N}\,d\nu=N$.  Moreover $\lambda_N\to\infty$ and
$b_{\lambda_N}\to c$ locally uniformly.  Let $Q_N$ be the canonical
$N$-particle gas with one-body weight $a_{\lambda_N}\nu$ and pair
interaction $\beta k_a$.  Then the original unweighted canonical gas
satisfies exactly
\[
 \frac{dP_{N,\beta}}{dQ_N}(\eta)
 =\frac{e^{-\frac\beta2 C_a(\eta-b_{\lambda_N}\nu,
                              \eta-b_{\lambda_N}\nu)}}
 {E_{Q_N}e^{-\frac\beta2 C_a(\eta-b_{\lambda_N}\nu,
                              \eta-b_{\lambda_N}\nu)}}.
\]
\end{corollary}
\begin{proof}
The angular formula in the proof of
Proposition~\ref{prop:saddle-reference} gives $C*b_\lambda\ge0$.
Hence \eqref{eq:reference-saddle} implies
$b_\lambda\le\lambda u$, where
$\int u\,d\nu=(p-1)^{-1}<\infty$.  The total mass is consequently
continuous in $\lambda$ by dominated convergence on compact activity
intervals, strictly increasing, and tends to zero as $\lambda\downarrow0$.
By monotone convergence and \eqref{eq:saddle-local-limit}, it tends to
$c\nu(\mathbb D)=\infty$ as $\lambda\to\infty$.  This proves the
assertions about $\lambda_N$ and the reference limit.  In
Theorem~\ref{thm:actual-reference-tilt}, choose
$\zeta_N=\lambda_Ne^{-\beta g(a)/2}$ and condition both sides on total
count $N$.  Every such sector has positive probability.  The left
conditional law is exactly $P_{N,\beta}$, and the right reference law
is $Q_N$, giving the displayed identity.
\end{proof}

\begin{remark}[What the small norm establishes, and what it does not]
The positive kernel
$T(z,dw)=a_\lambda(w)(1-e^{-\beta k_a(d(z,w))})\nu(dw)$
has a weighted row sum strictly below one in the stated range.
Consequently its repeated path integrals are summable.  More explicitly,
if the right side of \eqref{eq:reference-influence-bound} is $q<1$,
then the total weight of paths starting at $z$, of any positive length,
whose endpoint is at distance at least $L$ from $z$, is bounded by
$e^{-\alpha L}q/(1-q)$.  To verify this, insert the weight
$e^{\alpha\sum d(z_{i-1},z_i)}$, use the triangle inequality, and
integrate successive vertices; an $n$-step path has weighted integral at
most $q^n$.  This includes hyperbolic volume growth exactly in the row
integrals.  It is the quantitative smallness required by finite-range
comparison and disagreement constructions for the reference gas.

Equation \eqref{eq:actual-exact-tilt} still contains the smooth, long-range
quadratic energy.  A nonnegative energy can create dependencies after
normalization; its positivity alone does not allow the reference path
bound to be transferred to the tilted law.  A uniform estimate on this
tilt under local changes, including its normalizing ratios, remains
unproved.  Furthermore $b_\lambda\to c$ is a theorem about the explicitly
constructed reference; it is not an identification of the intensity of
$P_{\zeta,\beta}$ or of its fixed-$N$ subsequences.  One can work
directly with the canonical identity in
Corollary~\ref{cor:canonical-reference}, but conditioning the reference
on its total count introduces an additional global dependency which the
grandcanonical row-norm bound alone does not control.  Alternatively,
passing a result for the grandcanonical mixture back to every $N$
requires a comparison of nearby particle numbers.  The gain here is an exact
starting point for resummation with a proved activity bound and a small
finite-range remainder, rather than an assumed background density or a
replacement of the original interaction by a massive one.
\end{remark}

\subsection{An exact restoring potential in the imaginary-field representation}

The short-range expansion can be controlled uniformly even at complex
activities.  This provides a further rigorous step in a Debye construction.
The following elementary version includes its proof to specify which
uniformity is obtained.

\begin{proposition}[Complex-activity coercivity for a repulsive gas]
\label{prop:complex-coercivity}
Let $A(x)m(dx)$ be a finite nonnegative measure, let $v(x,y)\ge0$ be a
measurable symmetric pair interaction, and set
\[
 p(x,y)=1-e^{-v(x,y)},\qquad
 q=\sup_x\int A(y)p(x,y)m(dy).
\]
Define the complex partition function
\[
 \mathcal Z_\phi=
 \sum_{n\ge0}\frac1{n!}\int
       \prod_i A(x_i)e^{i\phi(x_i)}
       \prod_{i<j}e^{-v(x_i,x_j)}\,dm^n.
\]
If $q\le1/19$, then $\mathcal Z_\phi\ne0$ for every real measurable
$\phi$, and
\begin{equation}\label{eq:complex-coercivity}
 \log|\mathcal Z_\phi|-\log\mathcal Z_0
 \le-\frac12\int A(x)(1-\cos\phi(x))m(dx).
\end{equation}
\end{proposition}
\begin{proof}
Put $f=-p$.  For $n\ge2$ let
$U_n(\mathbf x)=\sum_{G\text{ connected}}\prod_{\{i,j\}\in E(G)}f(x_i,x_j)$,
and put $U_1=1$.  The repulsive tree-graph bound is
\begin{equation}\label{eq:repulsive-tree-bound}
 |U_n(\mathbf x)|\le
       \sum_{T\text{ tree}}\prod_{\{i,j\}\in E(T)}p(x_i,x_j).
\end{equation}
Here is a direct verification.  Root a connected labelled graph at vertex
$1$, form its distance levels, and give each other vertex the smallest
labelled parent among its neighbors in the preceding level.  This maps
the graph to a rooted tree $T$.  The graphs giving $T$ are precisely
$T\subset G\subset R(T)$, where $R(T)$ permits all edges within a level
and all edges to the preceding level whose endpoint there is at least
the chosen parent's label.  Summing over this interval of graphs gives
\[
 \prod_{e\in T}f_e\prod_{e\in R(T)\setminus T}(1+f_e).
\]
Since $0\le1+f_e\le1$, summing these disjoint intervals proves
\eqref{eq:repulsive-tree-bound}.

Fixing a root, successive integration of leaves bounds a tree integral
by $q^{n-1}$ times the root integral.  There are $n^{n-2}$ labelled
trees, as counted by their Pr\"ufer codes.  Consequently the connected
series
\begin{equation}\label{eq:connected-complex-series}
 L_\phi=\sum_{n\ge1}\frac1{n!}
       \int U_n(\mathbf x)\prod_i A(x_i)e^{i\phi(x_i)}\,dm^n
\end{equation}
converges absolutely when $eq<1$, uniformly in the phase.  More generally
it converges when the activity is multiplied by a complex scalar $t$
with $|t|eq<1$.  The identity $\exp L_\phi=\mathcal Z_\phi$ first
holds as a power-series identity near $t=0$: each graph decomposes
uniquely into connected components, which is precisely the exponential
formula for the partition series.  Analytic continuation within
$|t|eq<1$ proves it at $t=1$.  This also proves nonvanishing and
$\operatorname{Re}L_\phi=\log|\mathcal Z_\phi|$.

For any real numbers $\theta_1,\ldots,\theta_n$,
\[
 1-\cos(\theta_1+\cdots+\theta_n)
 \le n\sum_{j=1}^n(1-\cos\theta_j).
\]
Indeed, use the telescoping bound
$|\prod_j e^{i\theta_j}-1|\le\sum_j|e^{i\theta_j}-1|$
and Cauchy--Schwarz, then divide the squared inequality by $2$.
The $n=1$ term in $\operatorname{Re}(L_\phi-L_0)$ is
$-I$, where $I=\int A(1-\cos\phi)\,dm$.  The preceding inequality,
the tree bound and rooted leaf integration show that the sum of the
absolute values of all remaining terms is at most $S(q)I$, with
\[
 S(q)=\sum_{n\ge2}\frac{n^n}{n!}q^{n-1}
 \le\frac{e^2q}{1-eq}.
\]
The last bound uses $n!\ge(n/e)^n$.
For $q\le1/19$, the elementary bound $e<11/4$ gives
$S(q)\le121/260<1/2$.  Thus
$\operatorname{Re}(L_\phi-L_0)\le-(1-S(q))I\le-I/2$, proving
\eqref{eq:complex-coercivity}.
\end{proof}

Apply this proposition to the reference activity $A=a_\lambda$, measure
$m=\nu$, and pair interaction $v=\beta k_a$.  If
$0<\beta\le1/20$ and $a=\sqrt\beta$, then
$q\le\beta e^\beta\le1/19$, using
$e^x\le(1-x)^{-1}$ for $0\le x<1$. All bounds are uniform in
$\lambda$, including the choices $\lambda=\lambda_N$. This statement
concerns the grandcanonical partition functional at those activities,
not the partition function conditioned on count $N$. Indeed, a constant
phase $\phi_0$ multiplies a canonical partition function by
$e^{iN\phi_0}$ and leaves its modulus unchanged. Thus the canonical
constant phase is an explicit noncontracting mode.

There is an exact Gaussian formula behind this estimate.  Let $Y_a$ be
a centered real Gaussian field with covariance $C_a$.  It can be
constructed by sphere-averaging a zero-boundary Gaussian free field and
adding the independent harmonic Gaussian series with covariance $h$.
The averaged field has a measurable version, finite at every interior
point; its local mean-square continuity also follows from the integrable
logarithmic circle averages defining $C_a$.  The random variable
$b(Y_a)=\int b(z)Y_a(z)\nu(dz)$ is well defined in $L^2$, since
$C_a(b\nu,b\nu)<\infty$.  If $\Xi_{\rm short}$ denotes the
unnormalized grand partition functional of the finite-range pair
interaction, then
\begin{equation}\label{eq:exact-imaginary-field}
 \Xi_{\zeta,\beta}
 =e^{\frac\beta2C_a(b\nu,b\nu)}
 E\left[e^{-i\sqrt\beta b(Y_a)}
       \Xi_{\rm short}(a_\lambda e^{i\sqrt\beta Y_a})\right].
\end{equation}
To verify it, expand the finite grand partition series and take the
Gaussian characteristic function of
$\sqrt\beta(\eta-b\nu)(Y_a)$.  This gives exactly
$e^{-\frac\beta2C_a(\eta-b\nu,\eta-b\nu)}$, and the identity follows
from the square completion above.  Fubini is justified because the
absolute value before Gaussian averaging is bounded by the finite
positive partition function $\Xi_{\rm short}(a_\lambda)$.

For $\beta\le1/20$, Proposition~\ref{prop:complex-coercivity} proves the
pointwise field bound
\begin{equation}\label{eq:actual-field-restoring}
 \left|e^{-i\sqrt\beta b(Y_a)}
       \Xi_{\rm short}(a_\lambda e^{i\sqrt\beta Y_a})\right|
 \le\Xi_{\rm short}(a_\lambda)
       \exp\left[-\frac12\int a_\lambda(z)
             (1-\cos(\sqrt\beta Y_a(z)))\nu(dz)\right].
\end{equation}
On the part of the disk where $|\sqrt\beta Y_a|\le\pi$, the integrand
in the exponent is at least
$\beta a_\lambda Y_a^2/\pi^2$, since
$1-\cos t\ge2t^2/\pi^2$ for $|t|\le\pi$.
Thus the exact resummation has a quadratic restoring term in its small
field region.  It is periodic in the field and has additional wells.
There is also an imaginary linear residue: the derivative of
$\log\Xi_{\rm short}(a_\lambda e^{i\sqrt\beta Y})$
at $Y=0$ is $i\sqrt\beta\int\rho_{Q_\lambda}Y\,d\nu$.
After multiplication by the centering phase in
\eqref{eq:exact-imaginary-field}, it is
$i\sqrt\beta\int(\rho_{Q_\lambda}-b)Y\,d\nu$.
The constructed reference profile $b$ is not asserted to equal
$\rho_{Q_\lambda}$, so this term has not been canceled.
Controlling these wells, the boundary field and ratios of the complex
integrals in \eqref{eq:exact-imaginary-field} remains necessary.  An
upper bound on the absolute integrand does not give a lower bound on
its oscillatory expectation or a comparison of normalized local laws.
These are the remaining large-field and normalization steps; they are
not proved by the small reference norm or by
\eqref{eq:actual-field-restoring}.


\section{The reference-density residue and an exactly centered reference}
\label{sec:reference-residue}

The reference profile $b_\lambda$ in
Proposition~\ref{prop:saddle-reference} is a deterministic Poisson--Boltzmann
profile.  It is not the intensity of the finite-range reference gas.
We first bound their difference, and then construct a different reference
whose centering is exactly its own intensity.  These assertions concern
the reference representation; neither identifies the intensity of the
original disk gas.

Throughout this section, $0<\beta\le1/20$, $a=\sqrt\beta$,
$c=4/\beta-1$, $p=2-\beta/2$, $u(z)=(1-|z|^2)^p$,
$k=k_a$, and $M=M_a=2\log\cosh(a/2)$.  We use the abbreviation
\[
 t=\beta cM\le\beta(1-\beta/4)\le\frac{79}{1600}<\frac1{20}.
\]
For a finite nonnegative activity $A\nu$, let $Q_A$ denote the
grandcanonical gas with this activity and pair interaction $\beta k$,
and let $\rho_A$ be its intensity relative to $\nu$.

\begin{lemma}[Elementary insertion identities]
\label{lem:reference-insertion}
For every such activity,
\begin{align}
 \rho_A(z)&=A(z)E_{Q_A}\exp\!\left[-\beta\sum_{w\in\eta}k(d(z,w))\right],
 \label{eq:reference-gnz}\\
 \rho_A(z)&\le A(z),\qquad
 \rho_A^{(2)}(z,w)\le A(z)A(w).
 \label{eq:reference-ruelle}
\end{align}
If $(k*\rho_A)(z)$ is finite, then also
\begin{equation}\label{eq:reference-jensen}
 \rho_A(z)\ge A(z)\exp[-\beta(k*\rho_A)(z)].
\end{equation}
\end{lemma}
\begin{proof}
Expand the finite partition series and distinguish one of its labelled
points.  Reindexing the remaining count proves
\eqref{eq:reference-gnz}.  Distinguishing two points gives the analogous
formula with the extra factor $e^{-\beta k(d(z,w))}$ and both
insertion factors.  All these factors lie in $[0,1]$, which proves
\eqref{eq:reference-ruelle}.  Tonelli applies because the series are
nonnegative and bounded by the Poisson partition series.  Finally
Jensen's inequality applied to the exponential in
\eqref{eq:reference-gnz} proves \eqref{eq:reference-jensen}.
\end{proof}

\begin{proposition}[A uniform second-order density residue]
\label{prop:reference-density-residue}
Let $b=b_\lambda$ and $A=b\exp(\beta k*b)$ be the reference activity
in Theorem~\ref{thm:actual-reference-tilt}.  Then, uniformly in
$\lambda>0$ and $z\in\mathbb D$,
\begin{equation}\label{eq:reference-density-residue}
 |\rho_A(z)-b(z)|\le2\beta^2 b(z).
\end{equation}
\end{proposition}
\begin{proof}
Write $s(z)=\beta(k*b)(z)$, so that $0\le s\le t\le\beta$ and
$A\le e^t b$.  By the insertion identities and Jensen,
\[
 \rho_A\ge A e^{-\beta k*A}
 \ge b\exp[-(e^t-1)s]\ge b e^{-d},
 \qquad d=t(e^t-1)\le\beta^2e^\beta.
\]
Since $e^\beta<2$, this implies $\rho_A\ge(1-2\beta^2)b$.

We need a square-integral bound for the short kernel.  The decomposition
$k=q_a+S_aq_a$, Jensen's inequality for the averaging operator, and its
preservation of $\nu$ give
\[
 \int k(d(o,z))^2\nu(dz)\le4\int q_a(d(o,z))^2\nu(dz).
\]
Integration by parts in the normalized radial element
$\frac12\sinh r\,dr$ gives
\[
 \int q_a^2\,d\nu
 =\int_0^a q_a(r)\tanh(r/2)\,dr
 \le\frac12\int_0^a r\log(a/r)\,dr=\frac{a^2}{8}.
\]
Here $q_a(r)=\int_r^a ds/\sinh s\le\log(a/r)$; the boundary terms
vanish, including at zero since $r^2\log^2r\to0$.  Thus
\begin{equation}\label{eq:reference-k-square}
 \int k(d(o,z))^2\nu(dz)\le\frac{a^2}{2}.
\end{equation}

Fix $z$ and put $P_z(w)=1-e^{-\beta k(d(z,w))}$ and
$r_z=\frac{\beta^2}{2}(k^2*b)(z)$.  Since $b\le c$, the preceding
bound gives $0\le r_z\le\beta^2ca^2/4\le\beta^2$.
The inequalities $1-e^{-x}\ge x-x^2/2$ and $1-e^{-x}\le x$ give
\[
 \int P_z b\,d\nu\ge s(z)-r_z,
 \qquad \int P_z A\,d\nu\le e^t s(z).
\]
For any finite list $p_i\in[0,1]$, the even Bonferroni inequality is
$\prod_i(1-p_i)\le1-\sum_i p_i+\sum_{i<j}p_ip_j$.
For example, it follows by applying the two-term union bound to
independent events of probabilities $p_i$.  Apply this inequality to
the product in \eqref{eq:reference-gnz}; then use the lower bound
$\rho_A\ge e^{-d}b$ and the upper bound on the second correlation.
The result is
\[
 \frac{\rho_A(z)}{b(z)}
 \le e^{s(z)}\left[1-e^{-d}(s(z)-r_z)
                         +\frac12e^{2t}s(z)^2\right].
\]
All terms are integrable by the estimates already obtained.  Since
$e^s(1-s)\le1$ for $s\ge0$ and $1-e^{-d}\le d$, this yields
\[
 \frac{\rho_A(z)}{b(z)}
 \le1+e^t\left(td+r_z+\frac12e^{2t}t^2\right)
 \le1+\beta^2 e^\beta
                    \left(\beta e^\beta+1+\frac12e^{2\beta}\right).
\]
For $\beta\le1/20$, use $e^\beta\le20/19$ and
$e^{2\beta}\le10/9$ to bound the last coefficient by
$5500/3249<2$.  This proves the upper bound and hence the proposition.
\end{proof}

\subsection{Existence and uniqueness of an exact centering}

The residue can in fact be removed from the grandcanonical linear
field term.  The new centering is implicit; its construction does not
assume that the original gas is locally neutral.

\begin{theorem}[An exactly self-consistent finite-range reference]
\label{thm:exact-reference-centering}
For every $\lambda>0$ there is a unique radial function
$m_\lambda$ in the order interval $0\le m_\lambda\le\lambda u$
such that
\begin{equation}\label{eq:exact-reference-centering}
 m_\lambda=\rho_{A_\lambda},\qquad
 A_\lambda=\lambda u\exp[-\beta(C_a*m_\lambda)].
\end{equation}
Both functions are positive and continuous in the open disk and have
finite $\nu$-mass and finite $C_a$-energy.  They satisfy the uniform
activity estimate
\begin{equation}\label{eq:exact-centering-activity}
 0<m_\lambda\le A_\lambda
 \le c\exp(t e^{2t}).
\end{equation}
In particular, uniformly in $\lambda$,
\begin{equation}\label{eq:exact-centering-row}
 \sup_z\int A_\lambda(w)
       (1-e^{-\beta k(d(z,w))})\nu(dw)
 \le t e^{t e^{2t}}
 \le\frac{113918}{2180800}<\frac1{19}.
\end{equation}
With the additional weight $e^{d(z,w)}$ inside the integral, the
row sum is at most $e^{2\sqrt\beta}/19<2/19<1$.
The maps $\lambda\mapsto m_\lambda$ and
$\lambda\mapsto A_\lambda$ are continuous in $L^1(\nu)$ and in
the supremum norm.
\end{theorem}
\begin{proof}
Put $G_a=S_a^2g$, so that $C_a=G_a+h$.
The kernel $G_a$ is nonnegative, continuous, and positive semidefinite,
and its diagonal is the constant $g(a)$.  Positivity of the two by two
Gram matrices therefore gives $0\le G_a(z,w)\le g(a)$.
Continuity also follows directly by first averaging to obtain the
bounded continuous radial function $g(\max\{r,a\})$, and averaging
once more.  If $m$ is radial and integrable, angular averaging gives
$h*m=0$, with absolute integration justified by the boundedness of
$h(z,\cdot)$ for each fixed $z$.  Thus $C_a*m=G_a*m$.

Consider the convex set
\[
 \mathcal K_\lambda
 =\{m\in L^1(\nu):m\hbox{ is radial},\ 0\le m\le\lambda u\}.
\]
It is weakly compact.  To see this without a finite-volume assumption,
write $m=\lambda u f$ with $0\le f\le1$ in the finite measure
$u\nu$; its unit order interval is weak-star compact in $L^\infty(u\nu)$,
and integration against any bounded function gives its continuous image
in the weak topology of $L^1(\nu)$.
For $m\in\mathcal K_\lambda$ let
$A_m=\lambda u e^{-\beta G_a*m}$ and $T(m)=\rho_{A_m}$.
The insertion upper bound gives $0\le T(m)\le A_m\le\lambda u$,
and rotation covariance makes $T(m)$ radial.  Hence $T$ maps
$\mathcal K_\lambda$ into itself.

We verify continuity in the indicated weak topology.  Weak convergence
$m_j\to m$ implies $(G_a*m_j)(z)\to(G_a*m)(z)$ at every $z$,
since $G_a(z,\cdot)$ is bounded.  Dominated convergence then gives
$A_{m_j}\to A_m$ in $L^1(\nu)$.  Finite repulsive gases depend
continuously in total variation on their finite activity measures.
Indeed, telescope the $n$ activity factors in the partition series:
when their masses are bounded by $L$, the sum of the absolute differences
of the unnormalized sector densities is at most
$e^L\|A-B\|_1$.  The partition functions are at least one, so
normalization preserves convergence.  Insertion formula
\eqref{eq:reference-gnz}, whose expectation is of a function in $[0,1]$,
now gives
\[
 \|\rho_A-\rho_B\|_1
 \le\|A-B\|_1+\|B\|_1\operatorname{TV}(Q_A,Q_B).
\]
Thus $T(m_j)\to T(m)$ in $L^1$ and, in particular, weakly.
The Schauder--Tychonoff fixed-point theorem on a compact convex set
now gives a fixed point.  The preceding argument applies equally to
nets, or one can use that the finite measure $u\nu$ has separable
$L^1$ and its bounded weak-star order interval is metrizable.

Uniqueness follows from an entropy identity.  If $m,n$ are two fixed
points, the logarithm of the density ratio of their reference gases is
a constant minus $\beta\eta(G_a*(m-n))$.  This linear statistic is
integrable, because its summand is bounded and both expected counts are
finite.  Adding the two relative entropies therefore gives
\begin{align*}
 D(Q_{A_m}\Vert Q_{A_n})+D(Q_{A_n}\Vert Q_{A_m})
 &=\int(m-n)\log(A_m/A_n)\,d\nu\\
 &=-\beta G_a((m-n)\nu,(m-n)\nu)\le0.
\end{align*}
The energy is nonnegative by positive semidefiniteness and boundedness
of $G_a$, by finite-measure approximation.  Relative entropies are
nonnegative, so the reference laws coincide and their intensities
$m,n$ coincide.  This proves uniqueness.

The fixed point is positive by \eqref{eq:reference-gnz}: the empty
configuration has positive reference probability.  Its activity is
continuous by the continuity and boundedness of $G_a$ and dominated
convergence.  For each finite configuration,
$z\mapsto\exp[-\beta\sum_{w\in\eta}k(d(z,w))]$ is continuous,
with value zero at a coinciding point.  Here $k$ is continuous away
from zero and tends to infinity at zero.  Dominated convergence in
\eqref{eq:reference-gnz} proves continuity of $m$.
Finite mass follows from $m,A\le\lambda u$.  The boundary asymptotic
of $u$, with $p>1$, gives
$\int[-\log(1-|z|^2)]m\,d\nu<\infty$; the energy justification in
Theorem~\ref{thm:actual-reference-tilt} therefore applies.

We next prove parameter continuity, needed to select the bounded
branch in the uniform estimate.  If $\lambda_j\to\lambda>0$, the
profiles lie in one weakly compact order interval
$0\le m\le\Lambda u$.  Every weak limit is a fixed point at
$\lambda$, by the continuity just proved and the change in the scalar
activity.  Uniqueness identifies it with $m_\lambda$, and the same
argument gives strong $L^1$ convergence of the intensities and activities.
Furthermore
\[
 \|G_a*(m_j-m)\|_\infty\le g(a)\|m_j-m\|_1.
\]
Consequently the activities converge uniformly, and the pointwise
insertion formula, the uniform bound $A_j\le\Lambda$, and total
variation convergence of their laws give uniform convergence of the
intensities.  This proves both asserted continuities.  As
$\lambda\downarrow0$, $\|m_\lambda\|_\infty\le\lambda$.

It remains to prove the uniform bound.  Write $m=m_\lambda$, and set
\[
 w(z)=\log(\lambda u(z))-\beta(C*m)(z).
\]
For a radial nonnegative integrable charge, angular averaging yields
\[
 (C*m)(r)=\int-\log\max\{r,|y|\}\,m(y)\nu(dy)\ge0.
\]
It is finite and continuous and has boundary limit zero.  Writing
$B=C*m$, differentiation of that one-dimensional formula gives
\[
 B'(r)=-\frac2r\int_0^r\frac{s\,m(s)}{(1-s^2)^2}\,ds.
\]
Thus $B$ is twice continuously differentiable for $r>0$ and satisfies
$\Delta_{\mathbb H}B=-m/2$.  At zero both $B''(r)$ and $B'(r)/r$
tend to $-m(0)$; the radial Cartesian Hessian therefore extends
continuously there, proving the same classical equation at the center.
Since $\Delta_{\mathbb H}\log u=-p=-\beta c/2$,
\begin{equation}\label{eq:centering-w-equation}
 \Delta_{\mathbb H}w=\frac\beta2(m-c).
\end{equation}
The insertion Jensen bound and the fixed-point equation imply
\[
 m\ge A e^{-\beta k*m}=e^w.
\]
Also $w\to-\infty$ at the boundary, so $w$ attains an interior
maximum.  At that maximum \eqref{eq:centering-w-equation} gives
$m\le c$.  Hence $e^w\le c$ everywhere.  It follows that
\[
 m\le A=e^w e^{\beta k*m}
       \le c\exp(\beta M\|m\|_\infty).
\]
Put $x(\lambda)=\|m_\lambda\|_\infty/c$.  It is continuous,
tends to zero with $\lambda$, and satisfies $x\le e^{tx}$.
It cannot ever reach $2$, because $e^{2t}<2$ in the present range.
Thus $x<2$ for all activities and then $x\le e^{2t}$.  Substitution
in the preceding activity estimate proves
\eqref{eq:exact-centering-activity}.

Finally $1-e^{-\beta k}\le\beta k$ proves the first inequality in
\eqref{eq:exact-centering-row}.  For $0\le t\le79/1600$,
\[
 t e^{t e^{2t}}
 \le\frac{t}{1-t/(1-2t)}
 =\frac{t(1-2t)}{1-3t}
 \le\frac{113918}{2180800}<\frac1{19}.
\]
The first step uses $e^x\le(1-x)^{-1}$ twice.  The last rational
function is increasing on this interval, as direct differentiation
shows, and its value at $79/1600$ is the displayed fraction.
The weighted bound follows from the support condition $d(z,w)<2a$;
here $2\sqrt\beta<1/2$ and $e^{1/2}<2$.
This completes the proof.
\end{proof}

\begin{corollary}[The exactly centered representation, including counts]
\label{cor:exactly-centered-representation}
Set $\zeta=\lambda e^{-\beta g(a)/2}$, and use the reference of
Theorem~\ref{thm:exact-reference-centering}.  Then
\begin{equation}\label{eq:exactly-centered-tilt}
 \frac{dP_{\zeta,\beta}}{dQ_{A_\lambda}}(\eta)
 =\frac{\exp[-\frac\beta2 C_a(\eta-m_\lambda\nu,
                                  \eta-m_\lambda\nu)]}
 {E_{Q_{A_\lambda}}\exp[-\frac\beta2 C_a(\eta-m_\lambda\nu,
                                  \eta-m_\lambda\nu)]}.
\end{equation}
In the corresponding imaginary Gaussian representation, the derivative
at zero of the centered logarithmic reference functional vanishes:
for every bounded real test function $f$,
\begin{equation}\label{eq:exactly-centered-linear-cancel}
 \left.\frac d{ds}\right|_{s=0}
 \left[-is\sqrt\beta\int m_\lambda f\,d\nu
  +\log\Xi_{\rm short}(A_\lambda e^{is\sqrt\beta f})\right]=0.
\end{equation}
The periodic restoring estimate of
Proposition~\ref{prop:complex-coercivity} applies to this reference
uniformly in $\lambda$.
\end{corollary}
\begin{proof}
The exact fixed-point activity identity is
$\log(\lambda u)=\log A_\lambda+\beta C_a*m_\lambda$.
Use it in \eqref{eq:actual-smearing-density} and complete the square,
as in Theorem~\ref{thm:actual-reference-tilt}.  This proves the identity
in every count sector, including count zero.  The finite energies and
strictly positive denominator are justified exactly as there, using
$m_\lambda\le\lambda u$.  The logarithmic derivative of a grand
partition function is its intensity paired with the activity
perturbation.  Since $\rho_{A_\lambda}=m_\lambda$, the two terms
in \eqref{eq:exactly-centered-linear-cancel} cancel exactly.
All differentiations are justified by Poisson domination of the
finite partition series for bounded $f$.  The row bound
\eqref{eq:exact-centering-row} proves the final assertion.
\end{proof}

\begin{corollary}[Matching the reference mass to each canonical count]
\label{cor:exact-centering-mass}
The mass $\mathcal M(\lambda)=\int m_\lambda\,d\nu$ is continuous
and strictly increasing from zero to infinity.  Thus for every
$N\ge1$ there is a unique $\widehat\lambda_N$ for which
$\mathcal M(\widehat\lambda_N)=N$.  Conditioning
\eqref{eq:exactly-centered-tilt} at this parameter on count $N$
gives an exact representation of the original canonical gas.
\end{corollary}
\begin{proof}
Continuity follows from the $L^1$ continuity in the theorem, and
$\mathcal M(\lambda)\le\lambda\int u\,d\nu$ gives its zero limit.
For two parameters $\lambda>\kappa$, the same entropy calculation as
in the uniqueness proof gives
\begin{align*}
 &D(Q_{A_\lambda}\Vert Q_{A_\kappa})
  +D(Q_{A_\kappa}\Vert Q_{A_\lambda})
  +\beta C_a((m_\lambda-m_\kappa)\nu,
                    (m_\lambda-m_\kappa)\nu)\\
 &\hspace{12mm}=(\log\lambda-\log\kappa)
                    (\mathcal M(\lambda)-\mathcal M(\kappa)).
\end{align*}
Every term on the left is nonnegative.  Equality of the two masses
would imply equality of the reference laws and hence of their
intensities.  Their activities would also coincide, since the ratio
of the singleton-sector density to the vacuum probability equals the
activity density.  Their activity equations would then force
$\lambda=\kappa$.  Thus the masses are strictly increasing.
At the origin,
\[
 A_\lambda(0)=\lambda e^{-\beta(G_a*m_\lambda)(0)}
 \ge\lambda e^{-\beta g(a)\mathcal M(\lambda)}.
\]
The uniform activity upper bound implies
$\mathcal M(\lambda)\ge
 [\log\lambda-\log(c e^{t e^{2t}})]/[\beta g(a)]$,
so the mass diverges.  The intermediate value theorem gives the
unique matching parameter.  Every count sector has positive
probability, so conditioning the exact identity is legitimate.
\end{proof}

\begin{remark}[The gain and the remaining screening problem]
Exact centering removes one previously explicit obstruction: the
grandcanonical imaginary field action now has no linear term, while
its finite-range reference still has a uniform row norm below $1/19$.
The reference retains all particle-number weights.  This does not
remove the other periodic wells, prove a lower bound for oscillatory
normalizing integrals, or transfer reference mixing through the
long-range quadratic tilt.  It also does not assert
$m_\lambda\to c$, nor that the actual gas has intensity $m_\lambda$.
After conditioning the reference on count $N$, its intensity need not
equal $m_{\widehat\lambda_N}$; matching the expected total mass alone
does not preserve the grandcanonical linear cancellation for every
spatial test.  The canonical constant phase remains unchanged in
modulus.  These distinctions prevent the exact centering theorem
from being read as a proof of the intensity conjecture.
\end{remark}


\section{Exact self-consistent centering at fixed particle number}
\label{sec:canonical-self-centering}

The grandcanonical construction has a fixed-count counterpart. It removes
the linear field term at each finite $N$, but does not supply the uniform
reference estimates proved above for the grandcanonical gas.
Throughout this section $0<\beta<2$, $a>0$, $N\ge1$, and
$u(z)=(1-|z|^2)^p$, where $p=2-\beta/2>1$.

\begin{theorem}[Canonical centering]
\label{thm:canonical-self-centering}
There is a unique finite-energy self-consistent canonical reference
whose activity, up to a multiplicative constant, is
\[
 A_N(z)=u(z)\exp[-\beta(C_a*b_N)(z)]
\]
and whose intensity relative to $\nu$ is $b_N$.
Here the reference $Q_N^*$ has exactly $N$ particles and pair interaction
$\beta k_a$. The profile is radial and satisfies
\[
 \int b_N\,d\nu=N,\qquad
 e^{-\beta g(a)N}u\le A_N\le u.
\]
More precisely, finite energy in the uniqueness assertion means that
$b_N\nu$ has an absolutely defined Hilbert-space mean for the feature
map of $C_a$. The original canonical disk gas obeys the exact identity
\begin{equation}\label{eq:canonical-centered-exact}
 \frac{dP_{N,\beta}}{dQ_N^*}(\eta)
 =\frac{\exp[-\frac\beta2 C_a(\eta-b_N\nu,\eta-b_N\nu)]}
 {E_{Q_N^*}\exp[-\frac\beta2 C_a(\eta-b_N\nu,\eta-b_N\nu)]}.
\end{equation}
\end{theorem}
\begin{proof}
Since $C_a$ is a continuous positive semidefinite kernel on the disk,
it has a real separable Hilbert feature space $H$: take the span of
symbols $\Phi_z$, give them inner products
$\langle\Phi_z,\Phi_w\rangle=C_a(z,w)$, quotient by the null space,
and complete. Continuity of the kernel makes $z\mapsto\Phi_z$
continuous and a countable dense set of disk points spans a dense
subspace. In particular,
\[
 \|\Phi_z\|^2=D(z):=g(a)-\log(1-|z|^2).
\]
Let $Q_N^0$ be the canonical short-range gas with activity $u\nu$,
and set $X(\eta)=\sum_{z\in\eta}\Phi_z\in H$.
Its partition function is positive and finite because $k_a\ge0$ and
$\int u\,d\nu=(p-1)^{-1}$. For every $t>0$,
\begin{equation}\label{eq:canonical-feature-exp-moments}
 E_{Q_N^0}e^{t\|X\|}<\infty.
\end{equation}
Indeed $\|X\|\le\sum_z\sqrt{D(z)}$, so the unnormalized expectation
is bounded by a constant times
$[\int u(z)e^{t\sqrt{D(z)}}\nu(dz)]^N$. Writing
$s=-\log(1-|z|^2)$ evaluates the latter integral as
\[
 \int_0^\infty
 e^{-(p-1)s+t\sqrt{g(a)+s}}\,ds<\infty.
\]

Consider the functional on $H$
\[
 J_N(h)=\tfrac12\|h\|^2+\frac1\beta
       \log E_{Q_N^0}e^{-\beta\langle h,X\rangle}.
\]
The exponential moments justify its finiteness, norm continuity, and
Fr\'echet differentiation on every bounded set of $H$. For example,
the difference quotients and their first two derivatives are dominated
by a polynomial in $\|X\|$ times $e^{\beta R\|X\|}$ on
$\|h\|\le R$; these functions are integrable by
\eqref{eq:canonical-feature-exp-moments}.
If $Q_{N,h}$ is the probability measure obtained by the displayed
exponential tilt, then
\begin{align*}
 \nabla J_N(h)&=h-E_{Q_{N,h}}X,\\
 D^2J_N(h)[v,v]&=\|v\|^2+
       \beta\operatorname{Var}_{Q_{N,h}}\langle v,X\rangle
       \ge\|v\|^2.
\end{align*}
Thus $J_N$ is strongly convex. Jensen's inequality gives
\[
 J_N(h)\ge\tfrac12\|h\|^2-
               \langle h,E_{Q_N^0}X\rangle,
\]
so it is coercive. A minimizing sequence is bounded and therefore has
a weakly convergent subsequence in the Hilbert space. A norm-continuous
convex functional is weakly lower semicontinuous, so the weak limit is
a minimizer. Strong convexity makes this minimizer $h_N$ unique, and
its gradient equation reads
\begin{equation}\label{eq:canonical-feature-fixed-point}
 h_N=E_{Q_{N,h_N}}X.
\end{equation}

Put $Q_N^*=Q_{N,h_N}$ and let $b_N$ be its intensity density.
All first moments of $\sum\sqrt D$ remain finite under the tilt,
by the same exponential-moment argument. Consequently
\[
 h_N=\int\Phi_z b_N(z)\nu(dz),\qquad
 \langle\Phi_z,h_N\rangle=(C_a*b_N)(z),
\]
where the first integral is a Bochner integral and the second is
absolutely convergent by Cauchy--Schwarz.
This proves self-consistency. Conversely every self-consistent
reference in the stated finite-energy class defines a vector $h$
with $h=E_{Q_{N,h}}X$ and hence $\nabla J_N(h)=0$; uniqueness of
the minimizer gives uniqueness of both the reference and its intensity.

Rotations act unitarily on $H$ by $\Phi_z\mapsto\Phi_{e^{i\theta}z}$.
The kernel and $Q_N^0$ are rotation invariant, so $J_N$ is invariant
under these unitaries. Its unique minimizer is invariant, proving
radiality of the reference and its intensity.
For a radial background the harmonic part of $C_a$ has zero convolution,
and $C_a*b_N=G_a*b_N$, where $G_a=S_a^2g$.
The bounded positive semidefinite kernel $G_a$ satisfies
$0\le G_a(z,w)\le g(a)$: its diagonal is $g(a)$ and the upper bound
follows from Cauchy--Schwarz in its feature space.
These facts justify the displayed activity bounds because $\int b_N=N$.
The harmonic cancellation is legitimate even at the disk boundary:
for fixed $z$ its kernel is bounded as a function of $w$, and the
background has finite mass.

Finally, with $h=h_N$ and $X=X(\eta)$, completing the square gives
\[
 e^{-\frac\beta2\|X\|^2}\prod_{z\in\eta}u(z)
 =e^{\frac\beta2\|h\|^2}
  e^{-\frac\beta2\|X-h\|^2}
  \prod_{z\in\eta}u(z)e^{-\beta\langle\Phi_z,h\rangle}.
\]
Multiplying by the short-range interaction and using the exact
smearing identity for the disk gas proves
\eqref{eq:canonical-centered-exact}. Its denominator is strictly
positive and at most one, because the energy is a finite squared norm.
\end{proof}

\begin{corollary}[Cancellation of every canonical linear phase]
\label{cor:canonical-phase-centering}
Let $\mathcal Z_N(A)$ be the short-range canonical partition functional.
For every bounded real $\psi$, the function
\[
 F_N(t;\psi)=e^{-it\int b_N\psi\,d\nu}
               \frac{\mathcal Z_N(A_Ne^{it\psi})}
                    {\mathcal Z_N(A_N)}
\]
satisfies $F_N(0;\psi)=1$ and $F_N'(0;\psi)=0$.
If $W$ is an isonormal Gaussian process on $H$, the exact normalizer in
\eqref{eq:canonical-centered-exact} is
\begin{equation}\label{eq:canonical-gaussian-normalizer}
 E_W\left[e^{-i\sqrt\beta W(h_N)}
       \frac{\mathcal Z_N(A_Ne^{i\sqrt\beta W(\Phi_\cdot)})}
            {\mathcal Z_N(A_N)}\right].
\end{equation}
For a constant $\psi$, one has $F_N(t;\psi)=1$ for every real $t$.
\end{corollary}
\begin{proof}
The partition ratio is $E_{Q_N^*}e^{it\eta(\psi)}$; centering by
its mean gives the derivative assertions. A constant statistic equals
$N\psi$ deterministically, proving the last claim.
The Gaussian characteristic function gives
\[
 E_W e^{i\sqrt\beta W(X-h_N)}
       =e^{-\frac\beta2\|X-h_N\|^2}.
\]
Average this identity under $Q_N^*$ to obtain
\eqref{eq:canonical-gaussian-normalizer}. Fubini is justified because
the integrand before either expectation has modulus one. A jointly
measurable field $W(\Phi_z)$ can be used: expand in an orthonormal basis
of the separable feature space, take Gaussian partial sums, and choose
their measurable limit in local product $L^2$; its covariance is $C_a$.
Values on the resulting area-null exceptional set do not affect any
of the absolutely continuous partition integrals.
\end{proof}

The theorem does not provide bounds uniform in $N$ on short-range
influences or on the positive normalizer
\eqref{eq:canonical-gaussian-normalizer}. The grandcanonical insertion
identity used in the activity bootstrap compares the reference with
itself after insertion. At fixed count, insertion instead involves an
$(N-1)$-particle reference, whose intensity need not be $b_N$; the
same bootstrap therefore does not follow. Furthermore the constant
phase is exactly flat, so a canonical analogue of the grandcanonical
periodic coercivity estimate with a strictly positive penalty for all
constant phases is false. Exact centering removes the linear residue;
it does not remove these finite-volume normalization and comparison
problems.


\section{What has been removed from the boundary problem}
\label{sec:remaining-boundary-problem}

Theorems~\ref{thm:actual-tilt-TV}--\ref{thm:green-force-reconstruction}
are actual fixed-temperature estimates. They control all count and
position events for the specified perturbations, with constants
uniform in $N$. The force reconstruction also handles all distant
particles together. In particular, the Green force cannot retain
an extra random component at infinity for $0<\beta\le2$ after its
configuration-measurable finite truncations have been retained.
The independent radial theorem controls a further actual exterior
channel at every temperature.

At the determinantal benchmark, Theorem~\ref{thm:full-exterior-beta-two}
goes further: the entire far exterior loses influence on the whole
local law at rate $O_r(1-R^2)$, uniformly in $N$. Reversing the
conditioning reduces the estimate to the distant excess count
created by a fixed inner hole. Proposition~\ref{prop:general-hole-response}
formulates the same reduction for other temperatures under explicit
stochastic-domination and hole-response hypotheses.
The explicit discrete example in
Proposition~\ref{prop:lorentzian-no-covering} shows why Lorentzian
structure alone does not supply that domination. It is not a
counterexample for the disk interaction; the latter still has
additional structure which might yield a comparison.

These results leave two kinds of transmitted data unresolved. The
first is the harmonic contribution in the decomposition
$-\log|a-z|=g(a,z)-\log|1-a\bar z|$. The second is the scalar
normalization which governs local particle-number weights. Given
an exterior $Y$ and a prescribed count $m$, the exact conditional
density in a region $U$ is
\[
 \frac1{Z_{U,m,Y}}|\Delta_m(\mathbf z)|^\beta
       \exp\left\{\beta\sum_i\sum_{y\in Y}\log|z_i-y|\right\}
       \prod_i\mathbf1_U(z_i)\mu(dz_i).
\]
Adding a constant $c_0$ to the exterior potential has no effect in
one fixed sector, but multiplies its unnormalized weight by
$e^{\beta m c_0}$ when sectors are mixed. Force reconstruction
therefore does not determine the relative count probabilities.

The exact reference representation in
Section~\ref{sec:exact-smearing-reference} addresses part of this
normalization problem: it retains every count factor, constructs
the reference at every activity and at every canonical mass, and
proves a small short-range comparison norm without assuming the
actual intensity. It still requires an estimate on the normalized
long-range quadratic tilt under local changes. Dropping that tilt
because its energy is nonnegative would be invalid. Conditioning
the short-range reference on its total count also introduces a
dependency that the grandcanonical row-sum estimate does not
automatically control. For $\beta\le1/20$, the complex cluster
expansion additionally proves an explicit periodic restoring bound
for the exact imaginary-field integrand. The bound is quadratic
in each small-field well, but does not control other wells or
cancellation in the normalized oscillatory integral.

The new self-consistent construction in
Theorem~\ref{thm:exact-reference-centering} removes the imaginary
linear residue for $\beta\le1/20$ while preserving both the Mayer
bound and a uniformly subcritical distance-weighted influence norm.
Theorem~\ref{thm:canonical-self-centering} separately removes every
linear phase at each fixed $N$ for $0<\beta<2$. Its proof uses a
strongly convex Hilbert-space functional. The canonical theorem
does not give a reference influence estimate uniform in $N$, and
the constant canonical phase remains exactly flat. Neither
construction provides the normalized long-range comparison needed
for the original local particle law.

The profile $b_\lambda$ constructed there converges locally to
$4/\beta-1$. It is a deterministic balancing profile. It is not
asserted to be the intensity of either the interacting short-range
reference or the actual gas. Consequently the construction does
not prove
\[
 \mathbb E\Xi_\beta(dz)=(4/\beta-1)\nu(dz).
\]
The Green-force estimate likewise concerns a potential modulo
constants and cannot identify this coefficient by itself. These
are remaining estimates for the original model, not counterexamples
to its conjectured limit.

The susceptibility argument gives no comparable theorem for
$\beta>2$: the negative-type property used at its base is specific
to exponents at most two. Nor does the two-to-one coupling in the
appendix imply a coupling between every pair of successive particle
numbers. The simultaneous local linear response at $2$ gives an
additional perturbative input, while higher connected temperature
responses and a uniform remainder remain uncontrolled.

\appendix


\section{An exact non-determinantal transfer from two particles to one}
\label{sec:pair-deletion-below-two}

The next result restores a genuine count-transfer coupling below two.
Its scope is two interior particles versus one; the extension to
arbitrary interior counts is not established here.

\begin{proposition}[Deletion coupling for $0<\beta\le2$]
\label{prop:pair-deletion-below-two}
Let $\sigma$ be a finite nonatomic measure of compact support, let
$y\in\mathbb R^2$, and let $0<\beta\le2$.  Compare the unordered pair
law
\[
 P_2(dz\,dw)=\frac{|z-w|^\beta\sigma(dz)\sigma(dw)}
                    {\iint|u-v|^\beta\sigma(du)\sigma(dv)}
\]
with the singleton law
\[
 Q_y(dz)=\frac{|z-y|^\beta\sigma(dz)}
                   {\int|u-y|^\beta\sigma(du)}.
\]
There is a coupling $(\{Z,W\},X)$ with these marginals such that
$X\in\{Z,W\}$ almost surely.  In particular, for any finite exterior
$Y$ and observation region $U$, the actual disk-gas conditional law
with two particles in $U$ and exterior $Y$ dominates the conditional
law with one particle in $U$ and exterior $Y\cup\{y\}$.
\end{proposition}

\begin{proof}
Normalize $\sigma$ to have mass one and use the Hilbert embedding of
Lemma~\ref{lem:distance-negative-type}.  Put $h=\Phi(y)$, let $m$ and
$V>0$ be the mean and variance of $\Phi$ under $\sigma$.
For a Borel set $B$ with $b=\sigma(B)\in(0,1)$, put $a=1-b$.
Let $u,v$ be the conditional Hilbert means on $B,B^c$, and $S,T$ their
conditional variances.  Thus
\[
 m=bu+av,\qquad V=bS+aT+ab\|v-u\|^2.
\]
The probabilities of interest are
\[
 Q_y(B)=\frac{b(S+\|u-h\|^2)}{V+\|m-h\|^2},
 \qquad P_2(\{Z,W\}\subset B)=\frac{b^2S}{V}.
\]
We claim the first is at least the second.  Put $x=u-h$ and $d=v-u$.
The numerator after multiplying by the positive denominators and
cancelling $b$ is
\begin{align*}
 &V(S+\|x\|^2)-bS(V+\|x+ad\|^2)\\
 &\quad=a\left[
 b\bigl(\|d\|^2\|x\|^2-2S\langle x,d\rangle+S^2\bigr)
 +T\|x\|^2+aST\right]\ge0.
\end{align*}
The last inequality follows from Cauchy--Schwarz; the expression
in parentheses is at least $(\|d\|\|x\|-S)^2$.
Sets of $\sigma$-mass zero or one satisfy the same inequality directly.
Replacing $B$ by its complement gives, for every Borel $A$,
\begin{equation}
 Q_y(A)\le P_2(\{Z,W\}\cap A\ne\varnothing).
 \label{eq:pair-hall}
\end{equation}

For completeness, these Hall inequalities give the required continuum
coupling as follows.  Partition a compact set containing the support
into finitely many Borel cells of diameter at most $\varepsilon$.
The right states are unordered pairs of cells, including equal cells;
a left cell can be matched to a pair precisely when it belongs to that
pair.  Inequalities \eqref{eq:pair-hall}, applied to unions of cells,
are the finite capacitated Hall inequalities.  The max-flow/min-cut
theorem therefore supplies a coupling of the cell-valued singleton
and pair, supported on this incidence relation.  Lift the finite
coupling using the conditional laws within the cells, and within the
pair states.  It has the exact original marginals and satisfies
$\min(|X-Z|,|X-W|)\le\varepsilon$ almost surely.  Compactness gives a
weakly convergent subsequence as $\varepsilon\downarrow0$.
The marginals persist, and continuity of
$\min(|x-z|,|x-w|)$ forces its limiting value to be zero almost surely.
This is the desired deletion coupling.

For the disk-gas assertion, take
$\sigma(dz)=\mathbf1_U(z)\prod_{a\in Y}|z-a|^\beta dA(z)$.
The two displayed densities are exactly the two conditional laws.
\end{proof}

\begin{remark}
The deleted point has intensity measure $\rho_{P_2}-Q_y$, independently
of the choice of deletion coupling, because the difference consists
of one point.  This proves a count-sensitive transport statement for
the original interaction throughout $0<\beta\le2$.  It supplies no
distance decay by itself and does not justify iterating the coupling
through sectors with three or more particles.
\end{remark}


\section{A bounded response of the entire local law at temperature two}
\label{sec:whole-local-linear-response}

Here is an actual disk-gas estimate which includes every local count sector.
It concerns the derivative at $\beta=2$, and does not assume that a
general-temperature limiting process exists. Throughout this section,
$\mu=dA/\pi$, $W=B_r$ with $0<r<1$, and
\[
 S_N=\sum_{i<j}\log|z_i-z_j|,
 \qquad dP_{N,\beta}=Z_N(\beta)^{-1}e^{\beta S_N}\,d\mu^N.
\]
Let $P_{N,\beta}^W$ denote its law restricted to $W$. For a finite signed
measure $\sigma$, write $\|\sigma\|_{\mathrm{var}}=|\sigma|(\Omega)$;
the probability total variation distance is half this norm.

\begin{theorem}[Uniform whole-law linear response]
\label{thm:whole-law-linear-response}
For each $r<1$, the finite signed measures
\[
 \dot P_N^W=\left.\partial_\beta P_{N,\beta}^W\right|_{\beta=2}
\]
exist in variation norm, have total mass zero, and converge in variation
norm as $N\to\infty$. In particular, there is $C_r<\infty$ such that
\[
 \sup_N\|\dot P_N^W\|_{\mathrm{var}}\le C_r,
 \qquad
 |\operatorname{Cov}_{N,2}(F(\Xi_N|_W),S_N)|
       \le C_r\|F\|_\infty
\]
for every bounded measurable local observable $F$. The convergence is
simultaneous over $\|F\|_\infty\le1$, including all local count events.
\end{theorem}

This supplies a uniformly bounded local response to an extensive pair-energy
perturbation. It does not supply an interval of temperatures: a uniform
bound for a derivative at one temperature does not control its nonlinear
remainder uniformly in $N$.

\subsection{Kernel estimates with the collision singularities removed}
At $\beta=2$ the correlation kernel, relative to $\mu$, is
\[
 K_N(x,y)=\sum_{a=1}^N a(x\bar y)^{a-1},\qquad
 K(x,y)=(1-x\bar y)^{-2}.
\]
Write $\ell(x,y)=\log|x-y|$ and
$g_N(x)=\frac12\sum_{a=1}^N|x|^{2a}/a$,
$g(x)=-\frac12\log(1-|x|^2)$.

\begin{lemma}\label{lem:linear-log-kernel}
One has
\[
 \iint |\ell|\,|K_N-K|\,d\mu^2\longrightarrow0,
 \qquad
 J:=\sup_N\iint |\ell|\,|K_N|\,d\mu^2<\infty.
\]
Also $\int g\,d\mu=1/2$, and
\[
 L:=\sup_{z\in\mathbb D}\int_{\mathbb D}
            |\log|z-x||\,d\mu(x)\le\tfrac12+\log2.
\]
\end{lemma}
\begin{proof}
The last two assertions follow by integration: for the negative logarithm
translate into the unit disk about $z$, and for its positive part use
$|z-x|\le2$. Put $u=x\bar y$. The exact finite-series identity gives
\[
 |K_N-K|\le\frac{|u|^N}{|1-u|^2}
                   +\frac{N|u|^N}{|1-u|}.
\]
Where $\min(|x|,|y|)\le1/2$, denominators are bounded below and the
logarithmic singularity is uniformly integrable, so the integral tends to
zero. On the complement, put $v=|x||y|$ and let $\theta\in[-\pi,\pi]$
be the angular difference. Uniformly there,
\[
 |1-u|^2\ge c((1-v)^2+\theta^2),\qquad
 |\ell(x,y)|\le C(1+|\log|\theta||).
\]
Consequently the angular integrals with denominators $|1-u|^2$ and
$|1-u|$ are bounded respectively by
\[
 C\frac{1+|\log(1-v)|}{1-v},\qquad
 C(1+|\log(1-v)|^2).
\]
The radial integration identity
\[
 \int_0^1\!\int_0^1 f(rs)r\,dr\,s\,ds
       =\int_0^1 f(v)v(-\log v)\,dv
\]
shows that the first bound without $v^N$ is integrable, so its version
with $v^N$ tends to zero by dominated convergence. For the second bound,
split at $v=1/2$; on the upper half use $-\log v\le2(1-v)$ and
$v^N\le e^{-N(1-v)}$. Substituting $t=N(1-v)$ bounds the integral by
$C(1+\log^2(N+1))/N$, which tends to zero. The same integrable first
bound proves $\iint|\ell||K|<\infty$. The assertions follow.
\end{proof}

Fix $k\le N$ and distinct roots $\mathbf z=(z_1,\ldots,z_k)\in W^k$.
Set $G_N=[K_N(z_i,z_j)]$, $D_N=\det G_N$, and
\[
 Q_N(x,y)=K_N(x,\mathbf z)G_N^{-1}K_N(\mathbf z,y),
 \qquad q_N(x)=Q_N(x,x).
\]
These formulas are well defined: evaluations at $k$ distinct points are
linearly independent on the polynomials of degree less than $N$, by
polynomial interpolation. $Q_N$ is a rank-$k$ orthogonal projection,
and the reduced Palm kernel is $K_N-Q_N$. These statements also follow
directly by taking the Schur complement in the correlation determinant.

\begin{lemma}[Gram cofactor bounds]\label{lem:cofactor-linear-response}
Put $M=(1-r^2)^{-2}$ and $A=(1-r)^{-2}$. Uniformly in $N,k,\mathbf z$
as above and $x,y\in\mathbb D$,
\begin{align*}
 0\le D_N&\le M^k,\\
 D_N|Q_N(x,y)|&\le k^2M^{k-1}A^2,\\
 0\le D_N\{q_N(x)q_N(y)-|Q_N(x,y)|^2\}
       &\le4\binom{k}{2}^2M^{k-2}A^4.
\end{align*}
The last expression is zero when $k=1$. These bounds remain uniform when
the roots approach collisions.
\end{lemma}
\begin{proof}
The diagonal Gram entries are at most $M$ and the entries
$K_N(x,z_i)$ have absolute value at most $A$. For any two equally sized
index sets $I,J$, the Gram-minor Cauchy--Schwarz inequality gives
\[
 |\det(G_N)_{I,J}|^2
 \le\det(G_N)_{I,I}\det(G_N)_{J,J}\le M^{2|I|}.
\]
This inequality follows by writing the minor as the inner product of the
corresponding exterior products of Gram vectors. The first assertion is
Hadamard's inequality. Expanding $D_NQ_N$ using the adjugate of $G_N$
gives $k^2$ terms, each bounded by $M^{k-1}A^2$.

For the last assertion, form the $k$ by $2$ matrix $V$ of evaluations at
$x,y$. Its left side is $D_N\det(V^*G_N^{-1}V)$. Cauchy--Binet and
the complementary-minor formula for $G_N^{-1}$ express this as a sum
over pairs of two-element index sets. The complementary Gram minor is
bounded by $M^{k-2}$; each two by two evaluation determinant has absolute
value at most $2A^2$. There are $\binom{k}{2}^2$ terms. Positivity follows
because $V^*G_N^{-1}V$ is positive semidefinite.
\end{proof}

\subsection{The exact derivative and its all-order envelope}
Let $R_{N,\beta}^{(k)}$ be the factorial correlation density relative to
$\mu^k$. The angular mean identity for the logarithm gives
\begin{equation}\label{eq:response-h-potential}
 h_N(x):=\int\ell(x,y)K_N(y,y)\,d\mu(y)
   =-\frac12\sum_{a=1}^N\frac1a+g_N(x).
\end{equation}
Indeed, for a single term $a|y|^{2a-2}$, radial integration of
$\log\max(|x|,|y|)$ gives $(|x|^{2a}-1)/(2a)$.

Write $L_{\mathbf z}(x)=\sum_{i=1}^k\ell(z_i,x)$ and
$\Delta(\mathbf z)=\prod_{i<j}(z_i-z_j)$. Define
\begin{align}
 B_N(\mathbf z)={}&\sum_i g_N(z_i)-\int q_Ng_N\,d\mu
             -\int L_{\mathbf z}q_N\,d\mu\notag\\
 &+\frac12\iint\ell(x,y)
       \{q_N(x)q_N(y)-|Q_N(x,y)|^2\}\,d\mu^2\notag\\
 &+\iint\ell(x,y)\operatorname{Re}
       \{K_N(x,y)\overline{Q_N(x,y)}\}\,d\mu^2.
 \label{eq:whole-law-B}
\end{align}
Then the exact identity is
\begin{equation}\label{eq:whole-law-derivative}
 \dot R_N^{(k)}(\mathbf z):=
 \left.\partial_\beta R_{N,\beta}^{(k)}(\mathbf z)\right|_2
    =D_N\{\log|\Delta(\mathbf z)|+B_N(\mathbf z)\}.
\end{equation}
Here is a derivation to make the normalization cancellation explicit.
Differentiating the root-density integral gives
\[
 \frac{\dot R_N^{(k)}}{D_N}
  =\log|\Delta|+
   E^{!\mathbf z}_{N,2}\!\left[S_{N-k}+\sum_{y\in Y}L_{\mathbf z}(y)\right]
      -E_{N,2}S_N.
\]
Expand the one- and two-point densities of the Palm kernel $K_N-Q_N$.
The difference on the right, apart from $\log|\Delta|$, is
\begin{align*}
 \sum_i h_N(z_i)-\int q_Nh_N\,d\mu-\int L_{\mathbf z}q_N\,d\mu
 &+\tfrac12\iint\ell(q_Nq_N-|Q_N|^2)\,d\mu^2\\
 &+\iint\ell\operatorname{Re}(K_N\overline{Q_N})\,d\mu^2.
\end{align*}
Since $\int q_N\,d\mu=k$, the divergent harmonic-number constant in
\eqref{eq:response-h-potential} cancels exactly. This proves
\eqref{eq:whole-law-B}--\eqref{eq:whole-law-derivative}. All finite-$N$
differentiations are justified by integrability of logarithms times a
strictly positive Vandermonde power in a neighborhood of $\beta=2$.

\begin{proposition}\label{prop:response-all-order-envelope}
For each $r<1$ there are finite constants $c_r,C_r$ such that, for every
$k\ge1$ and $N\ge k$,
\[
 a_{N,k}:=\int_{W^k}|\dot R_N^{(k)}|\,d\mu^k
       \le c_r(k+1)^4C_r^k.
\]
For each fixed $k$, the functions $\dot R_N^{(k)}$ converge in
$L^1(W^k,\mu^k)$ as $N\to\infty$.
\end{proposition}
\begin{proof}
Lemmas \ref{lem:linear-log-kernel} and
\ref{lem:cofactor-linear-response}, applied term by term in
\eqref{eq:whole-law-B}, give
\[
 D_N|B_N|
 \le kM^kg(r)
  + k^2M^{k-1}A^2(\tfrac12+kL+J)
  +2\binom{k}{2}^2M^{k-2}A^4L.
\]
Call this upper bound $F_{r,k}$. If $v=\mu(W)=r^2$, then
\[
 a_{N,k}\le v^kF_{r,k}
       +M^k\binom{k}{2}v^{k-2}
          \iint_{W^2}|\ell(x,y)|\,d\mu^2.
\]
For $k=1$ the second term is zero, and for $k\ge2$ its double integral
is at most $vL$. This proves the claimed exponential times polynomial
bound after enlarging constants depending only on $r$.

For a fixed distinct root tuple, $G_N$ converges to its positive-definite
infinite Bergman Gram matrix. Positivity follows again by testing the
evaluation functionals against interpolation polynomials. Also
$K_N(x,z_i)\to K(x,z_i)$ uniformly for $|x|\le1$, so $Q_N\to Q$
uniformly on $(\overline{\mathbb D})^2$ for this tuple. Thus every term
in \eqref{eq:whole-law-B} converges: use $g_N\le g$, integrability of
the logarithms, and Lemma \ref{lem:linear-log-kernel} for the final
two-variable term. The bound just obtained, together with
$M^k\sum_{i<j}|\ell(z_i,z_j)|$, is an integrable pointwise majorant
on $W^k$. Dominated convergence proves $L^1$ convergence. The collision
set has $\mu^k$ measure zero.
\end{proof}

\subsection{Passing from every correlation to the entire local law}
\begin{proof}[Proof of Theorem \ref{thm:whole-law-linear-response}]
For finite $N$, differentiability of $P_{N,\beta}$ in variation norm
follows by differentiating its density in $L^1$. Its derivative is
$(S_N-E_{N,2}S_N)P_{N,2}$, and restriction to $W$ is a contraction for
variation norm.

For a second, more useful expression, let $j_{N,\beta,m}$ be the Janossy
density in $W$: the probability of exactly $m$ unordered points in $W$
is represented by $j_{N,\beta,m}\,d\mu^m/m!$. The empty sector uses
$\mu^0=1$. Finite inclusion--exclusion gives
\[
 j_{N,\beta,m}(\mathbf x)
 =\sum_{\ell=0}^{N-m}\frac{(-1)^\ell}{\ell!}
       \int_{W^\ell}R_{N,\beta}^{(m+\ell)}(\mathbf x,\mathbf y)
           \,d\mu^\ell(\mathbf y).
\]
Set $\dot R_N^{(k)}=0$ for $k>N$, and $\dot R_N^{(0)}=0$. Differentiate
this finite sum and use the triangle inequality. Regrouping by
$k=m+\ell$ yields
\begin{equation}\label{eq:response-Janossy-bound}
 \|\dot P_N^W\|_{\mathrm{var}}
  =\sum_{m\ge0}\frac1{m!}\int_{W^m}|\dot j_{N,m}|\,d\mu^m
  \le\sum_{k\ge1}\frac{2^k}{k!}a_{N,k}<C_r.
\end{equation}
The last inequality is uniform by Proposition
\ref{prop:response-all-order-envelope}.

Apply the same inequality to the difference of the derivatives for $N$
and $N'$. For every fixed $k$, the $L^1$ difference of the correlation
derivatives tends to zero. Its summand is dominated by twice the
summable envelope in \eqref{eq:response-Janossy-bound}. Hence
$\|\dot P_N^W-\dot P_{N'}^W\|_{\mathrm{var}}\to0$. Completeness of
the space of finite signed measures in variation norm gives the stated
limit. Its total mass is zero by continuity of total mass. Finally,
integrating $F$ against the derivative gives the covariance formula and
its uniform bound.
\end{proof}

\begin{remark}[What is still missing]
The uniform all-order bound above refers to correlation order $k$, not
temperature derivative order. It closes the passage from fixed
correlations to all local events at first order. A fixed interval near
$2$ still requires a uniform nonlinear remainder, or summable bounds on
all temperature derivatives. Neither follows from this theorem. For
example, an additional term proportional to $\tanh(N(\beta-2)^2)$ has
zero first derivative at $2$ for every $N$, while remaining nonzero at
every fixed other temperature as $N\to\infty$.
\end{remark}


\begin{thebibliography}{9}
\bibitem{BrandenHuh}
P.~Br\"and\'en and J.~Huh,
\emph{Lorentzian polynomials}, Annals of Mathematics
\textbf{192} (2020), 821--891.
\href{https://arxiv.org/pdf/1902.03719}{Original manuscript,
arXiv:1902.03719v8}; Theorems 2.25 and 2.30, Proposition 2.33.
The direct continuum susceptibility proof in this note is independent
of the polynomial characterization.
\bibitem{LyonsScreening}
R.~Lyons,
\emph{Determinantal probability: basic properties and conjectures},
Proceedings of the International Congress of Mathematicians 2014,
vol.~IV, 137--161.
\href{https://arxiv.org/abs/1406.2707}{arXiv:1406.2707}; Theorem~3.8.
The finite-projection comparison used here is also proved directly
in Lemma~\ref{lem:finite-projection-comparison}.
\end{thebibliography}
\end{document}
